% TaxoPolemics: A Formal Rhetoric of Disagreement
% Master dissertation file. This document carries the preamble, title page,
% and contents; the four papers are separate files, pulled in by \include
% (each \include starts its paper on a fresh page; \input would paste the
% file inline instead). Compile with pdflatex (run twice for the contents):
%   pdflatex dissertation.tex && pdflatex dissertation.tex
\documentclass[11pt]{report}

\usepackage[margin=1in]{geometry}
\usepackage{lmodern}
\usepackage{amsmath,amssymb}
\usepackage{graphicx}
\usepackage{booktabs,longtable,array,calc}
\usepackage{etoolbox}
\usepackage{xcolor}
\usepackage[colorlinks=true,linkcolor=black,urlcolor=blue]{hyperref}

% Pandoc-generated fragments use \tightlist in their lists.
\providecommand{\tightlist}{%
  \setlength{\itemsep}{0pt}\setlength{\parskip}{0pt}}

\title{TaxoPolemics\\[0.5em]\large A Formal Rhetoric of Disagreement}
\author{Jason Mark Wrinkle\\[0.3em]\normalsize\textit{with Yolanda Biche et al.}}
\date{November 21, 2026}

\begin{document}

\maketitle

\begin{abstract}
\noindent This dissertation presents TaxoPolemics, an engine for mapping
disagreement: a measured stance-space of topics and poles, an intensity
coordinate independent of leaning, and two formally distinct comparison
registers --- communion and conversation --- computed over the tensor
product of two profiles under a discipline that every headline number can
be un-contracted into its receipts. Chapter One traces the framework's
lineage from classical stasis theory through the modern theory of argument,
issue-based information systems, and knowledge organization. Chapter Two
states the methodology: research through design, taxonomy construction,
the measurement model, and the evaluation regime. Chapter Three is the
companion paper, \textit{How Maps Multiply}, which states the pair
mathematics of the founding domain. Chapter Four rehearses the defense:
the strongest anticipated objections, the responses, the concessions, and
the conditions under which the framework's claims would fall. The six pair
scores used throughout stand as provisionally ruled definitions
(2026-10-30) and are flagged as such wherever they appear.
\end{abstract}

\tableofcontents

\chapter{Lineage: The Ancestors of a Taxonomy of Argument}
\emph{Jason Mark Wrinkle (principal investigator and first author);
Yolanda Biche et al.~(drafting and computation). Working chapter,
2026-11-21. Status markers are kept visible: claims about the project's
own state are dated, and unsettled design questions are named as open
rather than smoothed over.}

\begin{center}\rule{0.5\linewidth}{0.5pt}\end{center}

\section{1.1 The claim this chapter
supports}\label{the-claim-this-chapter-supports}

TaxoPolemics is an engine for mapping disagreement. Given a domain of
difficult questions, it builds a \emph{hypermap}: a taxonomy of topics,
each topic articulated into two to eight \emph{poles} --- the stances
the topic admits --- and each person located in that space by rating,
per topic, how far they lean toward each pole and, independently, how
much the question matters to them. Two such profiles can then be
\emph{compared} by formal operations that distinguish two relations
ordinary language conflates: whether two people could \textbf{commune}
(the crossed register) and whether they could \textbf{converse} (the
straight register).

This chapter argues that the engine is not an orphan. It descends from
at least five identifiable traditions --- classical stasis theory, the
modern theory of argument, the design-methods tradition of issue-based
information systems, the study of classification itself, and the
psychometrics of attitude --- and that in each case the descent is
structural, not decorative: the ancestor solved a piece of the same
problem, and TaxoPolemics inherits both the solution and the problem it
left behind. A framework that can name its ancestors can also name its
delta. The final section does so.

\section{1.2 Stasis theory: the first taxonomy of
disagreement}\label{stasis-theory-the-first-taxonomy-of-disagreement}

The oldest ancestor is the most direct. In the second century BCE,
Hermagoras of Temnos founded what the rhetorical tradition calls
\textbf{stasis theory} (Greek \emph{stasis}, a stand, a position taken
in a dispute). Stasis theory holds that any controversy, to be argued at
all, must first be \emph{located}: the disputants must establish what
kind of question they are actually disagreeing about. The canonical
system distinguishes four stases --- \textbf{conjecture} (did it
happen?), \textbf{definition} (what is it?), \textbf{quality} (was it
right, how serious is it?), and \textbf{procedure} (is this the right
forum, the right process?). Quintilian and Cicero both transmit the
doctrine; it structured rhetorical education for over a millennium.

Stasis theory is, in modern vocabulary, a \emph{taxonomy of
disagreement} --- a claim that disputes have kinds, that the kinds can
be enumerated, and that identifying the kind is prior to, and more
valuable than, winning the argument. Every subsequent rhetorical
treatment of ``finding the real issue'' is downstream of it.

TaxoPolemics inherits the founding conviction --- that the first
intellectual act in a dispute is locative, not combative --- and changes
the unit of analysis. Stasis classifies \emph{disputes}; TaxoPolemics
classifies \emph{stances within topics}, and computes locations for
\emph{persons}. Where the stasis system asks ``what is this argument
about?'', the hypermap asks ``on this question, which poles exist, where
does each person stand among them, and how much does each person care?''
The ancient system is a diagnostic art exercised by a trained rhetor;
the modern one is an instrument any two disputants can run on
themselves. The continuity is the point: both rest on the premise that
most failed arguments fail because the parties never established what
they were disagreeing about, and that a shared map of the disagreement
is a precondition of a productive one.

\section{1.3 The modern theory of argument: Toulmin, Perelman,
Burke}\label{the-modern-theory-of-argument-toulmin-perelman-burke}

Three twentieth-century rhetoricians rebuilt the study of argument after
its long subordination to formal logic, and each contributes a
load-bearing element.

\textbf{Stephen Toulmin's} \emph{The Uses of Argument} (1958) proposed
that arguments have a \emph{layout} --- claim, grounds, warrant,
backing, qualifier, rebuttal --- and, more importantly for our purposes,
that the force of argumentative terms is \textbf{field-dependent}:
``probable,'' ``sufficient,'' ``reasonable'' mean different things in
law, in science, and at the dinner table, and no field-independent logic
governs them all. The hypermap system is field-dependent in exactly
Toulmin's sense, and wears the fact on its face: one engine runs sixteen
domains --- sexuality, politics, Abrahamic faith, chemistry, aesthetics,
philosophy, and others --- and in each domain the \emph{vocabulary of
the comparison itself} is re-skinned. The crossed relation is ``Play''
in the first domain and ``Commune'' in the vanilla domains; the straight
relation is ``Parlay'' in one and ``Converse'' in the others; the depth
measure is the Kink Factor in one and the Depth Factor elsewhere. A
field-independent framework would not need such skinning. TaxoPolemics
assumes Toulmin was right: the instrument must speak each field's own
language or it will not be believed in any of them.

\textbf{Chaïm Perelman and Lucie Olbrechts-Tyteca's} \emph{The New
Rhetoric} (1958) restored \emph{audience} to the centre of argument
theory: argumentation is the study of the means by which minds gain
adherence, and its basic data are not propositions but \emph{agreements}
--- the shared premises, values, and hierarchies from which a particular
audience can be moved. Their catalogue of the ``starting points of
argumentation'' is, like stasis, a taxonomy, and like this project, it
is interested less in who wins than in the prior question of what two
parties already hold in common. The straight register of the hypermap
--- the measurement of shared ground, in the same orientation, between
two profiles --- is Perelman's question rendered computable: not ``can
these two people persuade each other?'' but the antecedent ``what,
exactly, do they already agree on, and how strongly?''

\textbf{Kenneth Burke} contributes the project's central distinction. In
\emph{A Grammar of Motives} (1945) and the \emph{Rhetoric of Motives}
(1950), Burke argued that rhetoric is rooted in \textbf{identification}:
persuasion works not by defeating an opponent but by establishing
consubstantiality --- a sense of shared substance --- across division.
Identification is the correct ancestor for the crossed register's ruled
meaning on the vanilla maps: \textbf{communion}. The crossed comparison,
in which each person's leaning is read against the complementary leaning
of the other, asks whether two people can be \emph{of one substance} on
a topic --- can live it, work it, be it together --- as distinct from
the straight comparison, which asks only whether they would find each
other interesting to talk to. That the two relations are different
facts, priced separately by different mathematics (an indefinite form
for the crossed, a positive-definite one for the straight, as the
companion paper shows), is the project's central formal claim, and it is
a Burkean claim before it is a mathematical one.

\section{1.4 The design-methods tradition: IBIS and the wicked
problem}\label{the-design-methods-tradition-ibis-and-the-wicked-problem}

The second lineage is not rhetorical but computational, and it is the
closest technical ancestor. In the 1970s, Horst Rittel --- who had
already named the class of \textbf{wicked problems}, problems with no
definitive formulation and no stopping rule --- built, with Werner Kunz,
the \textbf{Issue-Based Information System (IBIS)}: a notation for
design deliberation with exactly three node types. An \textbf{Issue} is
a question under dispute; \textbf{Positions} are the stances one can
take on it; \textbf{Arguments} support or object to positions. IBIS was
implemented in software (notably Jeff Conklin's gIBIS in the 1980s) and
survives today in dialogue-mapping practice.

The resemblance to a hypermap topic is not accidental and should not be
understated: an IBIS Issue \emph{is} a topic, its Positions \emph{are}
poles, and the design intent --- externalise the structure of a
contested question so a group can deliberate over a shared map rather
than past each other --- is the same intent. TaxoPolemics is IBIS's
descendant, and the descent specifies the delta with unusual clarity:

\begin{enumerate}
\def\labelenumi{\arabic{enumi}.}
\tightlist
\item
  \textbf{Measurement.} IBIS positions are categorical: one adopts a
  position or one does not. A hypermap pole carries a signed magnitude
  --- a leaning, -100 to +100 --- so agreement and opposition become
  matters of degree, and two people who nominally hold ``the same
  position'' can be shown to hold it at different strengths and in
  different blends.
\item
  \textbf{The (n+1)th coordinate.} IBIS has no representation of
  \emph{how much an issue matters} to a participant, independent of the
  position taken on it. TaxoPolemics separates the leanings (n numbers
  for n poles) from \textbf{intensity} --- a signed relevance
  coordinate, the (n+1)th number --- as a matter of ruled design
  (2026-11-19). The separation is the project's single most
  consequential modelling decision, and \S 1.6 shows it has an independent
  ancestry in attitude research.
\item
  \textbf{Comparison as a first-class operation.} IBIS maps a
  deliberation; it does not compare two deliberators. The hypermap's
  pair mathematics --- the tensor product of profiles, its crossed and
  straight contractions, and the scores they yield --- has no IBIS
  analogue. IBIS tells a group what the questions are; TaxoPolemics
  additionally tells two people where they stand relative to each other,
  with receipts.
\item
  \textbf{Generativity of topics.} IBIS issues are authored in the room.
  Hypermap taxonomies are authored in advance, as \emph{skins} ---
  curated instruments, one per domain, sharing an engine. The
  deliberation is therefore instrumented rather than improvised: closer
  to a psychometric test than to a flip-chart, with the authoring burden
  (and its responsibilities, treated in Chapter Two) falling on the
  taxonomy's designer.
\end{enumerate}

\section{1.5 Knowledge organization: classification as
argument}\label{knowledge-organization-classification-as-argument}

The third lineage is the study of classification itself. Geoffrey Bowker
and Susan Leigh Star's \emph{Sorting Things Out: Classification and Its
Consequences} (1999) established, for information science, that
classification systems are neither neutral nor innocent: every taxonomy
encodes a theory of its domain, renders some distinctions visible and
others unthinkable, and has consequences for the people sorted by it.
The field of \textbf{knowledge organization} treats the design of a
classification --- its categories, its facets, its rules of combination
--- as scholarly work in its own right, judged by coherence, coverage,
and candour about its own choices.

A hypermap skin is a classification in exactly this sense, and the
project accepts the field's discipline. Two of its standing rules are
knowledge-organization rules. The \textbf{neutral umbrella rule} (ruled
2026-11-17) governs topics whose names join two stances with an
ampersand or slash: the joined name must be a \emph{neutral subject term
covering both sides}, never one of the joined terms itself --- a topic
may not smuggle a verdict into its own title. And the authoring standard
for topic descriptions (the \emph{hovertip}: the issue stated neutrally,
then one line per pole, each stated as its holders would state it) is a
rule of representational fairness of the kind Bowker and Star demand and
few production taxonomies observe. That a sexuality taxonomy, a politics
taxonomy, and a consciousness taxonomy (the sixteenth skin, Nouscapades,
added 2026-11-21) are built under one set of such rules is itself a
claim of the work: classification design is portable across domains
precisely when its fairness rules are stated at the level of the engine,
not improvised per domain.

\section{1.6 Attitude research: the independent ancestry of the (n+1)
model}\label{attitude-research-the-independent-ancestry-of-the-n1-model}

The separation of \emph{direction} from \emph{intensity} --- of which
way a person leans from how much the question matters to them --- was
arrived at in this project from instrument design (a slider's blend and
a slider's brightness are different facts, and conflating them corrupts
both). It is therefore worth recording that social psychology reached
the same separation from the opposite direction, empirically. The
\textbf{attitude strength} literature (Krosnick, Petty, and
collaborators, from the 1990s) established that an attitude's
\emph{direction and extremity} and its \emph{strength} --- importance,
certainty, accessibility --- are distinct properties with distinct
consequences: strong attitudes persist, resist persuasion, and guide
behaviour; weak ones, however extreme on paper, do not. A poll that
reports only leanings is therefore systematically misleading, in a way
the literature can quantify.

TaxoPolemics' (n+1) model is the attitude-strength distinction built
into the instrument rather than corrected for afterwards: every topic
rating is a vector of leanings \emph{plus} a signed intensity ---
positive for caring in favour of the issue's salience, negative for
active aversion to the issue itself --- displayed as brightness (full at
maximum positive, dimming to indifference at zero, passing into
complementary colours below zero). The companion tensor paper states the
pair mathematics for the two-flank case; the multi-pole generalisation,
including the open question of how the intensity coordinate should
weight pair scoring, is stated there as an open problem rather than
resolved by fiat. An honest lineage includes the note that the project's
most original modelling decision is also its least finished --- a state
of affairs Chapter Three's defence treats as a feature of the work's
stage, not an embarrassment to be hidden.

\section{1.7 Further kin: deliberative democracy and the map as
medium}\label{further-kin-deliberative-democracy-and-the-map-as-medium}

Two nearer relations deserve brief acknowledgement. \textbf{Deliberative
polling} (James Fishkin, from 1988) demonstrated that structured
exposure to balanced information changes what citizens think in
measurable, non-random ways --- evidence that disagreement is partly an
artefact of unstructured encounter, and that instruments can improve it.
The hypermap is a deliberative instrument in this tradition, though a
self-administered and pairwise one rather than a convened one. And the
project's commitment to the map as the primary object --- the ratings
are drawn as a navigable hypermap, not tabulated as a report ---
descends from the argument-mapping and information-visualisation
tradition (and from Tufte's insistence that a display is an argument), a
commitment Chapter Two treats methodologically: in this system, the
visualisation is not an illustration of the data but the form in which
the data is \emph{read}, and several of the project's defects (recorded
in the companion paper, \S 7) were precisely cases where a display quietly
stopped being the statistic it claimed to be.

\section{1.8 The delta, stated}\label{the-delta-stated}

Against these ancestors, the work claims four things as its own:

\begin{enumerate}
\def\labelenumi{\arabic{enumi}.}
\tightlist
\item
  \textbf{A measured stance-space.} Stasis located disputes; IBIS mapped
  them; neither measured the disputants. The hypermap locates
  \emph{persons} in a pole space with signed magnitudes, per topic, at
  scales from two to eight poles.
\item
  \textbf{Intensity as an independent coordinate.} The (n+1) model ---
  leanings plus a signed relevance --- as a ruled, first-class part of
  the rating, not a derived quantity.
\item
  \textbf{Two priced relations.} The crossed and straight comparisons
  --- communion and conversation --- as distinct mathematical objects
  (indefinite and positive-definite forms respectively) computed over
  the same product, where every predecessor system that scores
  ``compatibility'' at all produces a single number that confuses the
  two.
\item
  \textbf{Receipts.} Every headline number in the system can be
  \emph{un-contracted} --- opened back into the per-topic components,
  gates, and settings that produced it. A compatibility figure that
  cannot show its receipts is a horoscope; the project's founding
  discipline, stated in the companion paper and enforced throughout its
  development, is that no number ships without them.
\end{enumerate}

What the work does \emph{not} claim: that its taxonomies are complete,
that its scoring definitions are final (the six pair scores stand as
provisionally ruled definitions, 2026-10-30, pending review), or that
the instrument has been validated at population scale. Those are the
subjects of the next two chapters.

\begin{center}\rule{0.5\linewidth}{0.5pt}\end{center}

\emph{Next: Chapter Two --- Methodology: research through design,
taxonomy construction, the measurement model, and the evaluation
regime.}


\chapter{Methodology: Building an Instrument for Disagreement}
\emph{Jason Mark Wrinkle (principal investigator and first author);
Yolanda Biche et al.~(drafting and computation). Working chapter,
2026-11-21. The method described is the method as actually practised;
where practice falls short of the stated standard, the shortfall is
recorded rather than edited out.}

\begin{center}\rule{0.5\linewidth}{0.5pt}\end{center}

\section{2.1 Research through design}\label{research-through-design}

This work proceeds by \textbf{research through design}: the knowledge
claim is carried by an artefact --- the TaxoPolemics engine and its
sixteen domain taxonomies (``skins'') --- and by the design knowledge
extracted from building, breaking, and repairing it. The method is
appropriate to the question. Whether disagreement can be given a
measured, navigable form is not a proposition that can be settled in the
abstract; it is settled, if at all, by an instrument that either
sustains the form under real use across real domains or collapses under
it. The corpus of the research is therefore the shipped system: one
engine, sixteen skins (fifteen live at the time of writing; the
sixteenth, Nouscapades, committed to the repository 2026-11-21 with its
wiring deliberately held), roughly a thousand authored topics,
eighty-eight sample profiles, and a development record of several
hundred dated builds in which every significant design decision was
ruled, recorded, and --- frequently --- reversed on evidence.

Three consequences of the method are accepted at the outset. First, the
artefact is the primary text; this dissertation is its commentary, and
where the two disagree the artefact governs. Second, design decisions in
this project are made by \emph{ruling} --- a single investigator
decides, the decision and its date are recorded, and later evidence may
supersede it, with the supersession itself recorded rather than silently
edited away. The methodology thus produces an audit trail of judgement,
not merely of code. Third, evaluation (\S 2.6) is behavioural: claims
about the instrument are verified by observing it behave under the
conditions a real user runs, never by inspecting that the code
\emph{could} behave so.

\section{2.2 Taxonomy construction}\label{taxonomy-construction}

Each skin is authored, not generated. The construction procedure,
settled across the first fifteen skins and applied unchanged to the
sixteenth, is:

\begin{enumerate}
\def\labelenumi{\arabic{enumi}.}
\tightlist
\item
  \textbf{Domain delimitation.} The skin's territory is stated in one
  sentence and policed. (Precedent: Forescapades covers Abrahamic faith
  specifically, not ``religion'' in general; Nouscapades is restricted
  to the questions arising from Michael Levin's research programme in
  basal cognition, with the wider meaning-of-life canon explicitly
  excluded on the investigator's ruling, 2026-11-21, on the ground that
  the descriptive register of the source science would be corrupted by
  importing normative traditions wholesale.)
\item
  \textbf{Category structure.} Topics are grouped into categories (seven
  to eighteen per skin; subcategories where the domain's own structure
  supplies them, and only there). Depth is a property of the domain, not
  a target: the engine historically capped taxonomic depth at three
  levels, and the cap was lifted (2026-11-20) when the formalism was
  generalised to arbitrary depth.
\item
  \textbf{Topic articulation into poles.} Each topic is expressed as a
  question the domain genuinely contests, articulated into two to eight
  poles. Two-flank topics in the founding domain carry the complementary
  pair \emph{give/receive}; general topics carry two opposed stances (a
  slider), or three to eight stances (triangles, diamonds, and radial
  widgets). A topic that cannot be articulated into honest poles is not
  a topic and is removed.
\item
  \textbf{The hovertip standard.} Every topic carries a description in a
  fixed form: the \emph{issue}, stated neutrally; then one line per
  pole, each stated as that pole's holders would state it. This is the
  project's principal fairness instrument, applied to friend and foe
  positions alike.
\item
  \textbf{The neutral umbrella rule.} A topic name that joins two
  stances (by ampersand or slash) must be a neutral subject term
  covering both, never one of the joined terms. Names are audited
  against this rule in batches; violations are renamed.
\item
  \textbf{The meter word.} Each skin names what its intensity coordinate
  measures, in its own register --- \emph{interest} in the founding
  domain, \emph{stance} in politics, \emph{belief} in faith,
  \emph{taste} in aesthetics, \emph{conviction} in the philosophy of
  mind (ruled 2026-11-21), and so on. The word is part of the
  instrument's honesty: it tells the rater what quantity they are
  reporting.
\item
  \textbf{Archetype samples.} Each skin ships with authored sample
  profiles --- archetypal occupants of its stance-space, each rating
  every topic --- which serve simultaneously as calibration fixtures for
  the comparison mathematics, as worked examples for new users, and as
  the nearest thing the method has to a validation set: if the engine's
  archetype projection cannot recover a sample's intended archetype from
  its own ratings, the mathematics, not the sample, is presumed wrong.
  (This test caught a real defect: the projection once returned nine
  archetypes tied at 100\% for the founding domain's principal sample,
  because the underlying pair score discarded sign; the fix --- an
  identity-framed, same-flank similarity --- restored a discriminating
  ranking, afterwards confirmed by the investigator as matching his
  self-knowledge.)
\end{enumerate}

\section{2.3 The measurement model}\label{the-measurement-model}

A \textbf{profile} is a vector over the skin's topics. For a two-flank
topic it carries two signed values (give, receive), each in {[}-100,
+100{]}; for an n-pole topic it carries n signed leanings; and for every
topic it carries, as a ruled extension (2026-11-19), an independent
\textbf{intensity} coordinate --- the (n+1)th number --- signed, where
the sign denotes aversion to the issue itself rather than leaning within
it. Marks (hard limit, not interested, fantasy only, already tried)
attach to topics as flags with defined arithmetic roles.

Two disciplines govern everything computed from profiles:

\begin{itemize}
\tightlist
\item
  \textbf{Direction is never contracted.} Give and receive are
  complements, never flavours of one quantity; no statistic in the
  system adds them, averages them, or rolls them up into one another, at
  any level of any hierarchy. (The companion paper records the single
  explicit exception --- the Personal Kink Factor's per-topic maximum
  flank, adopted by ruling and flagged in the text as the sole permitted
  contraction over direction --- and records two historical violations,
  since repaired, in the founding domain's heat-map tooltips and
  aggregators. The discipline is stated as a law precisely because the
  system kept breaking it in small ways, and each break was a small lie
  told to a user about themselves.)
\item
  \textbf{Receipts.} Every aggregate is a \emph{contraction} of a
  per-topic table, and the table is always recoverable: any headline
  number can be un-contracted into the topics that voted, the gates they
  passed, and the settings in force. Statistics that failed this test in
  development were treated as defects regardless of their plausibility.
\end{itemize}

Pair comparison forms the \textbf{tensor product} of two profiles and
reads it through two pairing kernels --- the crossed kernel (Krossmatch,
$\bowtie$) and the straight kernel (Straight Match, $\odot$) --- under a system of
\textbf{gates} (hard limits remove components outright and are never
outvoted; a route table classifies per-topic outcomes; Heat, a
per-session threshold state, decides what counts tonight; fantasy-only
material contracts in the conversation register only; the Veil gates
what an observer is \emph{shown}, never what counts). The formalism, its
six provisional scores, and its worked examples are the subject of the
companion paper and are not re-derived here. What belongs to methodology
is the status discipline: the six score definitions stand as
\textbf{provisionally ruled} (2026-10-30) --- adopted for use, flagged
in every document that states them, and not hardened into claims until a
scheduled review that the press of other work has repeatedly deferred.
The dissertation reports them as provisional because that is what they
are.

\section{2.4 The map as instrument}\label{the-map-as-instrument}

The profile is \emph{read} as a map: topics drawn as nodes around a
central avatar, categories as branches, ratings as colour on
direction-split stems, intensity as brightness. This is not decoration.
The layout problem --- placing up to several hundred rated topics
legibly around a centre, minimising spoke crossings --- is solved by an
annealing optimiser (the solver; on the button, \emph{Anneal}) operating
under geometric laws stated in centre-avatar radii (branch throw 50
radii, subcategory segment 20, leaf nodes at three times their label-fit
width), with a per-phase quality gate (a flat tolerance of 2 crossings)
as the sole reject criterion. Layout quality is a methodological
variable, not an aesthetic one: a map whose crossings obscure the stems
is an instrument misreporting its data, and the project treats
unreadable layouts with the same seriousness as wrong numbers. The
solver's parameters, the gate, and the rejected alternatives (two
earlier ``tidy'' engines, ruled dead 2026-11-04 and not revived) are all
on the record.

\section{2.5 The comparison registers and their
semantics}\label{the-comparison-registers-and-their-semantics}

The crossed register answers: \emph{can these two people do this
together?} --- on the founding map, Play; on the vanilla maps, Commune.
The straight register answers: \emph{do these two people hold this in
common, in the same orientation?} --- Parlay on the founding map,
Converse on the vanilla maps, shipped between persons as Straight Talk.
The semantics are ruled, and one finding of the work is a
\emph{negative} result about them: for multi-pole topics, a code trace
(2026-11-19) established that the two registers have, in the current
implementation, \textbf{collapsed into one} --- the straight reading of
a pole topic substitutes the crossed numbers identically, pole topics
appear in neither register's match lists, and in group comparisons of
three or more their depth entries compute to zero. The investigator's
tentative distinction (crossed = communion; straight = an interesting
conversation) stands recorded as tentative. The dissertation claims the
\emph{distinction} as a contribution and the \emph{collapse} as an open
defect; conflating the two would be exactly the kind of overclaim the
receipts discipline exists to prevent.

\section{2.6 Evaluation}\label{evaluation}

The instrument is evaluated by four means, in descending order of the
weight the project gives them:

\begin{enumerate}
\def\labelenumi{\arabic{enumi}.}
\tightlist
\item
  \textbf{Behavioural verification.} Every claimed behaviour must be
  observed operating in the rendered interface, in a fresh profile,
  under the state a real user runs. Code inspection certifies mechanism,
  never behaviour; this rule was adopted after a defect class (a map
  shipped in a frozen state) passed an entire audit fleet whose checks
  were static. Automated fleets --- browser-driven test squads run
  against the packaged build, each load tagged so its traffic is
  excluded from the public counters --- extend the same principle: they
  report what the viewport did, not what the code says.
\item
  \textbf{Calibration fixtures.} The eighty-eight archetype samples
  (\S 2.2, step 7) function as a regression suite for meaning: known
  inputs, known intended outputs, run after every change to the
  comparison mathematics.
\item
  \textbf{Case replication across domains.} The strongest structural
  evidence the project offers is that one engine, unchanged in its
  mathematics, sustains sixteen independently authored taxonomies across
  domains as unlike as sexuality, chemistry, and consciousness --- and
  that domain-specific failures (a widget wrong for a one-door topic; a
  radar misplaced under the wrong subcategory) were detectable \emph{as}
  failures because the framework made their wrongness visible.
\item
  \textbf{Adoption signals.} Anonymous, always-on counters (a persistent
  random identifier per browser, the skin in use, gate passage and
  session events; survey \emph{content}, by contrast, is collected on an
  opt-out basis, and the distinction is stated in the privacy policy)
  give the work its only population-scale evidence to date: 262 unique
  visitors and 193 gate passages recorded by 2026-11-21. These numbers
  are reported here at their true weight --- evidence that strangers
  arrive and proceed, not evidence that the instrument improves their
  arguments. No controlled study of deliberative outcomes has been run;
  Chapter Three addresses the point directly.
\end{enumerate}

\section{2.7 Ethics and consent
architecture}\label{ethics-and-consent-architecture}

The founding domain made consent a design constraint before it was a
research-ethics heading, and the architecture generalises. A profile is
sovereign: its owner's hard limits remove material from every
computation and cannot be overridden by any majority in group readings.
The Veil ensures that what a person is \emph{shown} of another's map is
itself governed. Sync cards --- printed artefacts carrying a compressed
profile in QR codes, including a deliberately truncated biographical
field (975 characters) and terms of engagement (325) --- embody the same
principle physically: a person hands over exactly the profile they
chose, in a form that cannot be silently expanded. The investigator's
own profile, and his partner's, are the founding samples; the
dissertation follows the project's rule that no real person's ratings
are published as data beyond what they themselves placed in the sample
set.

\section{2.8 Limitations of the method}\label{limitations-of-the-method}

Stated plainly, because Chapter Three will be asked: the taxonomies are
single-author instruments, and their pole articulations encode one
investigator's judgement of where the honest disagreements lie --- the
hovertip standard disciplines the \emph{statement} of positions but
cannot neutralise the \emph{selection} of topics. The scoring
definitions are provisional. The sample corpus is authored, not drawn.
The audience is, at the time of writing, one paying member and some
hundreds of visitors. And the development method --- a single
investigator ruling by recorded fiat, assisted by AI collaborators whose
contributions are credited in the artefact's own About text ---
concentrates both the coherence and the blind spots of the work in one
place. The methodology does not hide these facts; it is, in large part,
a machine for making them visible.

\begin{center}\rule{0.5\linewidth}{0.5pt}\end{center}

\emph{Next: Chapter Three --- The Defense: anticipated objections and
responses.}


\chapter{How Maps Multiply: The Tensor Product of Profiles}
% Options for packages loaded elsewhere
\PassOptionsToPackage{unicode}{hyperref}
\PassOptionsToPackage{hyphens}{url}
\documentclass[
]{article}
\usepackage{xcolor}
\usepackage[margin=1in]{geometry}
\usepackage{amsmath,amssymb}
\setcounter{secnumdepth}{-\maxdimen} % remove section numbering
\usepackage{iftex}
\ifPDFTeX
  \usepackage[T1]{fontenc}
  \usepackage[utf8]{inputenc}
  \usepackage{textcomp} % provide euro and other symbols
\else % if luatex or xetex
  \usepackage{unicode-math} % this also loads fontspec
  \defaultfontfeatures{Scale=MatchLowercase}
  \defaultfontfeatures[\rmfamily]{Ligatures=TeX,Scale=1}
\fi
\usepackage{lmodern}
\ifPDFTeX\else
  % xetex/luatex font selection
\fi
% Use upquote if available, for straight quotes in verbatim environments
\IfFileExists{upquote.sty}{\usepackage{upquote}}{}
\IfFileExists{microtype.sty}{% use microtype if available
  \usepackage[]{microtype}
  \UseMicrotypeSet[protrusion]{basicmath} % disable protrusion for tt fonts
}{}
\makeatletter
\@ifundefined{KOMAClassName}{% if non-KOMA class
  \IfFileExists{parskip.sty}{%
    \usepackage{parskip}
  }{% else
    \setlength{\parindent}{0pt}
    \setlength{\parskip}{6pt plus 2pt minus 1pt}}
}{% if KOMA class
  \KOMAoptions{parskip=half}}
\makeatother
\usepackage{longtable,booktabs,array}
\newcounter{none} % for unnumbered tables
\usepackage{calc} % for calculating minipage widths
% Correct order of tables after \paragraph or \subparagraph
\usepackage{etoolbox}
\makeatletter
\patchcmd\longtable{\par}{\if@noskipsec\mbox{}\fi\par}{}{}
\makeatother
% Allow footnotes in longtable head/foot
\IfFileExists{footnotehyper.sty}{\usepackage{footnotehyper}}{\usepackage{footnote}}
\makesavenoteenv{longtable}
\setlength{\emergencystretch}{3em} % prevent overfull lines
\providecommand{\tightlist}{%
  \setlength{\itemsep}{0pt}\setlength{\parskip}{0pt}}
\usepackage{bookmark}
\IfFileExists{xurl.sty}{\usepackage{xurl}}{} % add URL line breaks if available
\urlstyle{same}
\hypersetup{
  hidelinks,
  pdfcreator={LaTeX via pandoc}}

\title{How Maps Multiply}
\usepackage{etoolbox}
\makeatletter
\providecommand{\subtitle}[1]{% add subtitle to \maketitle
  \apptocmd{\@title}{\par {\large #1 \par}}{}{}
}
\makeatother
\subtitle{The Tensor Product of Profiles}
\author{}
\date{Revision of 2026-11-21}

\begin{document}
\maketitle

{
\setcounter{tocdepth}{3}
\tableofcontents
}
\begin{quote}
\textbf{Revision note --- 2026-11-21.} This revision adds two sections
and changes nothing in \S \S 1--10 or in Appendix A. \textbf{\S 11} records
the multi-pole extension as it actually stands in the shipped engine:
topics of one to eight poles, the agreement score and Depth Factor by
which pole topics are currently compared, and three findings from the
code trace of 2026-11-19 --- for pole topics the crossed and straight
registers presently coincide, the Depth Factor's intensity is
sign-blind, and pole topics do not yet appear in either register's match
lists --- all recorded as open, with the owner's tentative distinction
(crossed as communion, straight as an interesting conversation) marked
tentative, not ruled. \textbf{\S 12} states the owner's ruling of
2026-11-19 that intensity is an independent (n+1)th coordinate of a
topic rating --- relevance, signed, negative for aversion to the issue
itself --- displayed as brightness; it records that the shipped
implementation has decoupled the drawing only, that the stored-(n+1)
build is specified and queued, and that the coordinate's weight in pair
scoring is an open question reserved to the owner. A currency note: the
formalism of this paper is the founding domain's, but the same engine
now runs fifteen live skins (a sixteenth, Nouscapades, was committed to
the repository today, wiring held), reading the same algebra under
re-skinned names --- Depth Factor for Kink Factor, Commune for Play,
Converse for Parlay. {[}Census corrected in this revision on the owner's
ruling, 2026-11-21: the paper's original count of 403 topics predated
the phantom-topic cleanup and the removals that followed it. The live
catalog --- sexcapades.json as served in build 2026-11-21\_0022, 18
categories and 44 subcategories --- holds \textbf{372 topics}, counted
structurally from the file the app loads, and every denominator in the
paper now reads accordingly: 372 topics, 744 components. (An earlier
draft of this note cited 381, from a unique-id census of the code
bundle; the structural count of the live taxonomy file governs.){]} The
Adam Peasant $\bowtie$ Eve Vale recomputation flagged on 2026-11-06 remains
outstanding and is not touched.
\end{quote}

\begin{quote}
\textbf{Revision note --- 2026-11-06.} This revision updates the paper
to Eve Vale's new 246-rating profile (she was a 10-rating fixture) and
renames Rowan to Adam Peasant throughout. Eve's middle name `Morgan' is
struck and deprecated --- she is now Eve Vale. His full name is now Adam
Peasant (formerly Rowan; the Rowan $\rightarrow$ Adam rename stands, and the surname
Peasant is now added). Adam Peasant's ratings are unchanged from Rowan's
(164 ratings), so the Adam Peasant $\bowtie$ Riley headlines are unchanged. The
Adam Peasant $\bowtie$ Eve Vale shared count is now 135 (was 10) and the
Play-considered count is 134. The Conversation-gate count, the Play
match count, and the weighted Play Kink Factor and Play Match Score for
Adam Peasant $\bowtie$ Eve Vale have NOT been recomputed --- the pair engine
must be re-run on the new profiles --- and the old-Eve Vale figures are
recorded in the flag in the worked headlines, not as current values.
Nothing else in the argument, conclusions, or authorship has changed.
\end{quote}

\emph{The Tensor Product of Profiles}

\emph{How the Sexcapades sync algorithm works, in the language it
secretly speaks. Revision working draft, 2026-10-29, from the 2026-10-28
revision. Prepared for the Samhain review, 2026-10-31.}

\textbf{Authors:} Jason Mark Wrinkle --- principal investigator and
first author; Yolanda Biche et al.~--- drafting and computation.

\emph{Revision note (working draft, not paper text): this revision,
worked 2026-10-31, incorporates the 2026-10-30 definitions and namings,
including the six scores as provisionally ruled for consideration. It
does not harden interim, provisional, or recommended items into rulings.
Items marked {[}Flag{]} or {[}Session{]} in this draft are not settled
here, and the six-score definitions in \S 4 are provisional until the
review after Saturday. Hierarchical weights and kinkiness composition
are throughout the contractions, stated at rule level; the worked
headlines in \S 4 remain hierarchical-only, because no kinkiness fixture
for those topics is on record and none is invented here. Personal and
Charted formulas are stated at the level the 2026-10-30 core closure
supports; open details from that closure are listed in \S 4 and in the
companion notes, marked as recommendations or as open, not as rulings.
Appendix A remains HELD for the full review sometime after Saturday.}

\subsection{Definitions index}\label{definitions-index}

\textbf{Profile.} One person's vector: for each of the 372 topics, a
give value and a receive value, each from -100 to +100, with marks
bolted on (hard limit, not interested, fantasy only, already tried). 744
components.

\textbf{Direction complements.} Give and receive are complements, never
two flavors of one quantity. Nothing in this paper ever adds a person's
give and receive together. Direction is never contracted into one number
anywhere in the system, with one explicit exception stated in \S 4: the
Personal Kink Factor takes the maximum flank per topic, signed, as its
per-topic direction reading. That maximum is the sole permitted
contraction over direction, by explicit switch, and it does not license
any other.

\textbf{Tensor product ($\otimes$).} The full grid of pairwise combinations of
two profiles, every component of one laid against every component of the
other. The product stage: everything, before anything decides what
counts.

\textbf{Krossmatch ($\bowtie$).} The product restricted to matched components:
each topic paired with itself, each direction paired with its
complement. The operation of the sync between two people. {[}Session:
wording to confirm --- Krossmatch as outer product plus swap kernel,
reading the crossed diagonal.{]}

\textbf{Complement map ($\kappa$).} Re-reads a profile in the crossed frame:
($\kappa$Q) at topic t, direction d, equals Q at t with the direction swapped.
The Krossmatch is the first profile laid component-wise against the
complement of the second.

\textbf{Swap ($\sigma$).} The direction exchange: give exchanges with receive,
and applying it twice returns home ($\sigma^2$ is the identity).

\textbf{Straight Match ($\odot$).} The product restricted to the parallel
sector: each topic paired with itself, each direction paired with
itself. The object of peer likeness. {[}Session: wording to confirm ---
straight as the same-direction diagonal, inner-product-like.{]}

\textbf{Straight Talk.} The shipped reading of the Straight Match
between two people: per-topic straight sums, ranked, banded (\S 9).

\textbf{Contraction.} Summing an index away. A contracted index is a
question answered in aggregate; a free index is an axis the reader still
gets to see. The product stage builds the outer product; the contraction
stage reads a diagonal and sums it --- the inner-product-like step.
Outer and inner are different stages, not different names for one thing.
{[}Session: wording to confirm.{]}

\textbf{Overlay.} Co-display: both profiles shown together, nothing
summed. Overlay answers no question; contraction answers one. The
component table (\S 4) is the zero-contraction reading; an overlay is its
visual cousin.

\textbf{Gates.} What may count, decided before any sum: hard limits
(which remove components entirely), the route table (which decides what
a component counts as), Heat (tonight's threshold), fantasy rules
(Conversation only), and Veil (which gates what is seen, never what
counts).

\textbf{Play Match Score (Play MS).} The Play outer number, the outer
ring on a bowtie tab: the geometric-mean combination stated in \S 4, over
Play-considered topics, fantasy excluded from numerator and denominator.
{[}Flag: provisional definition, 2026-10-30, for consideration.{]}
Formerly called the Match Score in this draft; the Play MS name governs.

\textbf{Personal Kink Factor (Personal KF).} The individual profile's
Kink Factor, shown as the inner ring on the profile tab. It is a
whole-map contraction, stated in \S 4: per topic, the maximum flank
(signed, proportional, epsilon-thresholded) plus the flag sum once,
hierarchically weighted and kinkiness-composed, additive in shape, over
all 372 topics, not clamped at 100. {[}Flag: provisional definition,
2026-10-30, for consideration.{]}

\textbf{Charted.} The individual profile's outer ring on the profile
tab. It is a coverage and attended share, stated in \S 4: per topic, once,
binary rated (epsilon-thresholded, either flank) plus the flag sum once,
hierarchically weighted and kinkiness-composed, over all 372 topics, no
clamp. Charted is the name used for this ring throughout this revision.
{[}Flag: provisional definition, 2026-10-30, for consideration.{]}

\textbf{Play Kink Factor (Play KF).} The Play inner number, the inner
ring on a bowtie tab: the geometric-mean combination stated in \S 4, over
shared non-limit-blocked topics. {[}Flag: provisional definition,
2026-10-30, for consideration.{]} The name Play KF supersedes Party KF,
which is retired and is not used in this revision.

\textbf{Parlay Kink Factor (Parlay KF).} The straight inner number, the
inner ring on a circledot tab, mirroring Play KF on the Straight Match.
{[}Flag: provisional definition, 2026-10-30, for consideration.
``Parlay'' alone means nothing; the full name is always Parlay Kink
Factor.{]}

\textbf{Parlay Match Score (Parlay MS).} The straight outer number, the
outer ring on a circledot tab, mirroring Play MS on the Straight Match.
The name Parlay Match Score is ruled, superseding the provisional
working name Straight Match Score. Its definition is provisional,
2026-10-30, for consideration. {[}Flag: provisional definition.{]}

\textbf{Hot Topics.} The partial contraction: sum over people and
channels, leave the topic index free --- a ranked list per topic (\S 4).

\textbf{Lit sums (S, C).} S = gg + rr, the straight sum, whose
contraction is resemblance; C = gr + rg, the crossed sum, whose
contraction is fit (\S 10).

\textbf{Dark sums (K, $\epsilon$).} K = gg - rr, the channel imbalance of a
resemblance, unchanged when partners are exchanged; $\epsilon$ = gr - rg, the
orientation of action-flow, which changes sign when partners are
exchanged (\S 10).

\textbf{Heat as a phase.} A pairing-level angle, declared per session,
never stored in a profile; only the difference of the two angles enters
any score (\S 10, provisional in part: whether arousal is a third phase or
a gain on Heat is open).

\textbf{Initiative.} Bookkeeping at the resolution layer, per live
channel, not a contraction; a candidate, not built (\S 10, provisional).

\textbf{Hierarchical weights.} Each of the 18 categories carries an
equal share; within a category, its direct topics and its subcategories
split that share equally, direct topics being peers of subcategories;
within a subcategory, its topics split the subcategory's share equally;
a starburst splits among its topics. Weights apply to sums across
topics, never to per-topic terms. The solver's energy is not weighted.

\textbf{Kinkiness.} One minus popularity, where popularity is
prevalence, not acceptability. The effective weight of a topic in a
calculation is its hierarchical share multiplied by its kinkiness,
across all calculations, with the solver's energy excluded unless the
owner says otherwise. Popularity itself is an aggregate over seven
sources, weighted by statistical strength, stated in \S 4. {[}Flag: the
strength numbers per source are not ruled, and the aggregated score's
use as the Popularity List sort key is unverified in implementation;
this paper states the rule, not a verified implementation.{]}

\subsection{1. A profile is a vector}\label{a-profile-is-a-vector}

Take one person. For every topic in the taxonomy --- 372 of them ---
they carry two numbers: how much they want to \emph{give} it, and how
much they want to \emph{receive} it, each on a scale from -100 to +100.
A profile is therefore a list of 744 numbers: a vector in a
744-dimensional space, one component per (topic, direction) pair. Bolted
onto some components are marks --- \emph{hard limit}, \emph{not
interested}, \emph{fantasy only}, \emph{already tried} --- flags that
will matter later, when the vector is used.

Nothing about two people has been said yet. A single profile just sits
there, the way a single vector does: complete, static, and entirely
about one point of view.

One terminological commitment, made now because everything turns on it:
the two directions are \emph{complements}, not two flavors of one
quantity. Nothing in what follows ever adds a person's give and receive
together. The algebra keeps them separate at every stage, and the one
operation that would blend them is the one operation the system refuses
to perform. Direction is never contracted, in a profile, in a pair, or
in a headline.

\subsection{2. The product: every pairing at
once}\label{the-product-every-pairing-at-once}

Comparing two people means forming their \textbf{tensor product},
written with the circled cross: Adam Peasant $\otimes$ Eve Vale. The tensor
product of two vectors is not a single number and not a single new
vector --- it is the full grid of pairwise combinations, every component
of one profile laid against every component of the other.

Most of that grid is meaningless (his rating for rope against her rating
for latex says nothing). The sync lives on one thin, meaningful slice of
it: for each topic, the two \emph{live channels} --- his give against
her receive, and his receive against her give. That slice is a named
object in its own right: \textbf{the Krossmatch}, written with the
bowtie --- Adam Peasant $\bowtie$ Eve Vale --- the product restricted to matched
components, where a topic is paired with itself and a direction paired
with its complement, never with itself. On each channel sits a raw
component, the product of the two sides' values. The Krossmatch is the
complete grid of what could flow between the two people, before anything
decides what counts. (\S 8 examines the object itself.)

Components then resolve into outcomes. The strength of a pairing is set
by the \textbf{weaker side}: his 100 against her 51 is a pairing of
strength 51, because a scene happens at the level both people actually
reach. Each live channel carries its own intensity, on the /100 scale of
the ratings, per flank --- the give-to-receive channel and the
receive-to-give channel are never averaged into one number, because
direction is never contracted. Each topic resolves, through a fixed
route table over the two channels' signed values, into one of a small
set of outcomes: Limit-blocked, Invisible, Match, Curious fit, Conflict,
One-sided, Apart.

This table of per-topic outcomes is the Krossmatch's component list,
resolved --- the un-summed, fully detailed truth of the pair. Everything
else the sync reports is a summary of this table, obtained by
\emph{contracting} it. Laid out without summing, the same table is an
overlay: co-display, both profiles shown together, answering no
question. Contraction is what turns the table into an answer.

\subsection{3. The gates: why the contraction is not a pure
sum}\label{the-gates-why-the-contraction-is-not-a-pure-sum}

In Einstein's convention, contraction means summing over an index and
letting it disappear. A naive contraction of the Krossmatch --- just
count the matches, divide by something, done --- would be dishonest in
specific, nameable ways, so the sum is \textbf{gated} before it is
taken:

\begin{itemize}
\tightlist
\item
  \textbf{Hard limits remove Krossmatch components entirely.} A limited
  topic never enters any numerator. It is not averaged away, not
  outvoted --- in a group it cannot be carried by majority, either. It
  stands in a visible, named list, outside the arithmetic. For a single
  person's own rings, a limit-only topic is excluded altogether from
  that person's calculation, in both rings, while the denominator stays
  at 372. The scored value therefore drops while the denominator does
  not shrink. That inconsistency is intentional and accepted: a hard
  limit carries a penalty in the headline, and the paper does not
  ``fix'' it by shrinking the denominator. A limit plus another mark is
  not covered by this rule; the limit-only rule is not extended
  automatically.
\item
  \textbf{The route table decides what a component counts \emph{as}.} A
  curious fit is not a small match; it is a different outcome class, and
  which classes enter which sum is a design decision, stated in advance.
\item
  \textbf{Heat sets tonight's threshold.} Heat is a state, not a trait
  --- a slider, bipolar around Neutral, set per person per session and
  never written into the profile. Raising it lets weaker and
  wider-gapped components count; lowering it below Neutral sheds curious
  fits first. Same tensor, different contraction, because the
  \emph{occasion} changed.
\item
  \textbf{Fantasy-only components contract only in Conversation mode.}
  In Play mode they are excluded; in Conversation they are fuel. Already
  Tried never changes a component's value at all --- it annotates it,
  and in the Kink Factor tables it adds as a flag, as \S 4 states.
\item
  \textbf{Veil is not a gate on the sum. It is a gate on the observer.}
  Each owner's Veil projects the Krossmatch onto what that owner is
  shown: their own veiled territory hides, the partner's shows
  deliberately as a teaser. Heat decides what counts; Veil decides
  what's seen.
\end{itemize}

\subsection{4. The contractions: how one object is read, in pairs and
alone}\label{the-contractions-how-one-object-is-read-in-pairs-and-alone}

\textbf{The component table itself} is the zero-contraction reading:
every topic, its outcome, its strength, the gap between the two sides
displayed rather than penalised. This is what two people actually talk
through.

\textbf{The six numbers, provisional (2026-10-30).} The six scores are
three pairs: Personal KF and Charted for one person, Play KF and Play MS
for a crossed pair, Parlay KF and Parlay MS for a straight pair.
{[}Flag: every definition in this paragraph is provisional, ruled
provisionally for consideration, until the review after Saturday. Final
inclusion tables are deferred.{]} All six are hierarchy-share $\times$
kinkiness weighted, once, numerator and denominator alike. Kinkiness is
one minus popularity (\S 4, below). Combination across members is the
geometric mean. Normalization is by the group maximum, pinned
2026-10-30: a group-max-zero contribution contributes zero and stays in
the denominator.

Ground and fit decide what counts. Marks plus fit existence determine
ground; ground is multiplied by the type of fit, positive or negative.
Per member, binary fit existence plus the table sum gives the Match
Score shape, and rated-or-not plus the table sum gives the Kink Factor
shape, each normalized by the group maximum and combined by geometric
mean over members. Fit sign sorts ground into common ground (positive)
or debate ground (negative); both are conversational fodder. {[}Flag:
sign divergence is provisional, not settled. The maximum rule --- the
maximum of a negative and a positive fit is the positive; negative only
if both directions are negative --- is provisional.{]} Debate topics
count in the Parlay MS numerator and denominator, like Shop talk and
On-the-fence topics. Bands display sections only; none of the four pair
numbers refers to banded topics. Peer-ground is dropped; it is not a
term in these definitions.

Epsilon is the fit and rating threshold. Epsilon = $e_0$ $\times$
e\^{}(-heat/100), in profile units, with no rescaling; $e_0$ = 1. At $e_0$ = 1
the range is about 0.37 to 2.72 over heat -100 to +100. A binary fit
holds where the absolute match term exceeds epsilon. Group heat is the
arithmetic mean of the individual heats at formation. For a single
profile, epsilon thresholds a rating, and Charted counts a topic where
the absolute rating on either flank exceeds epsilon, with epsilon from
that profile's personal heat. {[}Flag: the epsilon formula and shape are
under review, not settled; the current shape, $e_0$ = 1, stays in the app.
An $e_0$ = 10 variant was ordered, stopped, and is not reintroduced without
a new ruling.{]}

Inclusion tables carry positive entries only, no null or zero entries,
and no banded topics. The Play table is, for now, copied from the
Charted table; its final form is deferred. The Parlay table is: Favorite
1, Tried 1, Fantasy 1 (full 1). Teased does not enter the Parlay table
--- it is absent, not zero. Favorite was restored to the Parlay table as
a +1 factor, not zero. {[}Flag: provisional, final inclusion tables
deferred.{]}

\textbf{Play MS} is the geometric mean over members of the normalized
(binary fit + table sum), where binary fit per direction holds iff the
absolute Krossmatch term exceeds epsilon, normalized by the group
maximum. Play-considered topics are shared, not limit-blocked, and not
fantasy --- fantasy is excluded from the numerator \emph{and} the
denominator, so Play asks only about what could actually happen. One
scalar --- but it is never allowed to travel alone; its weights and
counts ride with it, because a percentage with its indices hidden is
exactly how summaries lie. The Krossmatch term is give $\times$ receive
crossed, normed /100.

\textbf{Play KF} is the maximum fit from the Krossmatch over its two
directions, times the geometric mean over members of (binary ``rated or
not'' + table sum), normalized by the group maximum, over shared
non-limit-blocked topics as denominator.

\textbf{Parlay KF and Parlay MS} are the same shapes on the Straight
Match, with the Parlay table and the Parlay considered set: shared
non-limit-blocked topics, limit-blocked out, Fantasy included in the
Parlay considered set, unlike Play. The Straight Match term is give $\times$
give, receive $\times$ receive, normed /100. Parlay MS mirrors Play MS; Parlay
KF mirrors Play KF.

\textbf{Play MS against Play KF.} The two stand in a subset relation, at
the same Heat: every Play match is counted in Play KF's ground, and
every Play-considered topic is Play-KF-considered, because
Play-considered topics are the shared, non-limit-blocked topics minus
fantasy. The reverse does not hold: fantasy ground counts for Play KF
and never for Play MS. Heat moves both, but through different channels,
through epsilon. Same tensor, same Heat, different questions --- and the
subset relation is what keeps the two answers honest with each other.

\textbf{Hot Topics} is a \emph{partial} contraction: sum over the people
and the Krossmatch's channels, but leave the topic index free. What
remains is a vector indexed by topic --- a ranked list of what this
group, tonight, has the most breadth and strength to talk about. A topic
limited by anyone is excluded from the ranking outright.

\textbf{Group sync} adds one more gathering rule: members count equally
and topics count by their hierarchical weight, and a topic carries only
on a strict majority --- more than half, ties do not count --- while
limits remain personal and un-votable, and anyone may exit. Everything
personal is sovereign; everything collective is majority. The group's
Kink Factor is a Play Kink Factor, gathered under that majority rule,
normalized by the group maximum; a group-max-zero contribution
contributes zero and stays in the denominator. In table terms, all four
pair numbers read the same way: for pairs, hierarchical shares; for
groups, majority counting. That table rule is stated here as a table
rule; it is not, by itself, a claim about any other surface. {[}Flag:
group definitions provisional, as above.{]}

\textbf{Weights.} Every contraction in this section is hierarchically
weighted, and kinkiness-composed. Each of the 18 categories carries an
equal share; within a category, its direct topics and its subcategories
split that share equally, direct topics being peers of subcategories;
within a subcategory, its topics split the subcategory's share equally;
a starburst splits among its topics. The weights apply to sums
\emph{across} topics, never to per-topic terms: a topic's own fit
intensity, its straight sum, and its crossed sum are computed exactly as
before, and only their gathering into a headline is weighted. Aggregates
are weighted; per-topic examples are unweighted. That distinction is
what keeps the worked arithmetic of \S 8 and \S 10 standing exactly as
written while the headlines move. The solver's energy is not weighted,
and kinkiness does not enter it unless the owner says otherwise.

\textbf{Popularity and kinkiness.} Kinkiness is one minus popularity,
and popularity is prevalence --- how common a thing is --- not
acceptability. The popularity aggregate is weighted by statistical
strength over seven sources: the Joyal survey of 2015, with 1,516
adults, and the Lehmiller survey of 2018, with 4,175 adults, as anchors
at factor 1.0; FetLife counts, discounted at factor 0.5 and
log-normalized; Reddit counts, discounted, pending; community ratings as
one merged discounted source; the project's own survey; and literature
priors, meaning the written-erotica primary bank plus semantic priors,
provisional and weakest. Log scaling weight is ln(size) $\times$ factor; the
prior weight is 0.5 $\times$ ln(5) $\approx$ 0.80, on the log scale. Aggregation is
ignore-missing: it runs over only the sources with data for a topic,
renormalized, with missing not counted as zero and a prior-only
fallback. Automatic submissions when a FetLife import fails (Type 1) are
excluded from this aggregate entirely. Literature priors are
provisional, qualitative estimates, never cited as counts, distinct from
the two published surveys and not double-counting them. The aggregated
score is also the Popularity List sort key as a rule; whether the
shipped sort key implements it is unverified and is flagged, not
asserted. {[}Flag: Reddit counts pending; literature priors
provisional.{]}

\textbf{Worked headlines, hierarchical-only.} Adam Peasant $\bowtie$ Riley share
73 topics. Under Conversation gates, 45 count: the flat ratio reads
62\%, and the hierarchically weighted Play Kink Factor reads 58\% ---
lower, because the counted topics sit disproportionately in
lighter-weight branches. Under Play gates, 15 match of 73 considered:
flat 21\%, weighted Play Match Score 16\%. Adam Peasant $\bowtie$ Eve Vale share
135 topics. Under Play gates, fantasy sets one topic aside, leaving 134
considered, with no topics limit-blocked. {[}Flag: the Conversation-gate
count, the Play match count, and the weighted Play Kink Factor and Play
Match Score for this pair have not been recomputed for Eve Vale's new
246-rating profile. They were computed on her old 10-rating profile ---
Conversation 9 count, flat 90\%, weighted Play Kink Factor 96\%; Play 7
matches of 9 considered, flat 78\%, weighted Play Match Score 82\% ---
and those figures are retained here only as a record, not as current
values. Recomputing them requires running the pair engine on the new
profiles. The direction of the weighted-vs-flat comparison for the new
Eve Vale is not stated, because it depends on which topics the new
counted set falls in.{]} The weights do not change any topic's own
outcome; they change how much each outcome weighs in the headline. These
headlines are hierarchical-only, and they predate the provisional
six-score definitions above; they are not kinkiness-weighted
recomputations, because no kinkiness fixture for those topics is on
record, and none is invented here. A kinkiness-weighted headline for
these pairs would require that fixture first. {[}Flag: these worked
headlines have not been recomputed under the provisional six-score
definitions.{]}

\textbf{A single profile, read two ways.} A profile also contracts
alone, with no second profile taking part. Its inner ring is the
Personal Kink Factor, and its outer ring is Charted. Both are whole-map
contractions, and both use the same contributions table, stated here at
the level the core closure supports.

The basis is the whole map: one person's whole-map positive share in
either direction, hierarchically weighted and kinkiness-composed. The
denominator is all 372 topics. Personal counts topics, once, by the
maximum flank; Charted counts topics, once, by binary rated inclusion.
The shape is additive expanse, like Charted, not a flank share
multiplied by an outcome weight.

For each topic, the Personal contribution is the maximum flank (signed,
proportional, epsilon-thresholded) plus the flag sum once: favorite 1,
tried +1, fantasy +0.5, curious 0, teased 0. Into-it is the rating
condition, not a flag. An unrated-but-marked topic has base 0 while its
flags still count; a limit-only topic counts nothing. A negative maximum
can pull a topic's contribution negative. Personal is not clamped at
100. {[}Flag: provisional definition, 2026-10-30, for consideration. The
maximum-flank rule is the explicit per-topic direction-contraction
exception stated in the definitions index.{]}

Charted is coverage, an attended share. Per topic, once: binary rated,
epsilon-thresholded on either flank, plus the flag sum once, with the
same table. Negatives are taken as absolute value to the binary: Charted
stays positive, and a negative counts as much as a positive. A
NotInterested mark alone counts as attended for Charted and adds nothing
to Personal. A limit plus another mark is excluded altogether, the limit
winning. The flag sum counts once per topic, not per flank. Charted is
not clamped.

Open details from the core closure are carried here as open, not as
settled text. With no recommendation yet: group Charted, where the
record is individual only; and whether the shipped code counts the flag
sum per topic or per flank, subject to an editor check this paper does
not assert.

\textbf{The rings, in one line.} The in-tab Conversation and Play switch
is removed; mode follows the tab. On a profile tab, the outer ring is
Charted and the inner ring is the Personal Kink Factor. On a bowtie tab,
the outer ring is Play Match Score and the inner ring is Play Kink
Factor. On a circledot tab, there are two rings, arranged to mirror the
bowtie: the outer ring is Parlay Match Score and the inner ring is
Parlay Kink Factor, both provisional in definition. In the header, the
Charted ring takes the former header-ring seat.

\subsection{5. What the language buys, and where it
stops}\label{what-the-language-buys-and-where-it-stops}

The tensor language is not decoration. It enforces the design's central
honesty: \textbf{every headline number can be un-contracted.}
Un-contracting always lands in the same place: the Krossmatch, component
by component --- the grid of flows is where the receipts live. A score
is a sum, and a sum can always be opened back into its components ---
the topics that voted, the gates they passed, the setting stamped on the
result. A compatibility figure that cannot show its receipts is a
horoscope; one that can is a measurement.

Where the analogy honestly stops: there is no geometry here. Topics are
not coordinates on a smooth manifold; there is no distance between
``rope'' and ``latex,'' no metric tensor in Einstein's sense, nothing
that transforms covariantly. This is algebra, not differential geometry
--- a product, a set of gates, and a family of contractions. But that is
enough. One object, honestly built, read in pairs and alone, with no
reading allowed to hide how it was summed: that is the whole algorithm,
and it fits in one line ---

\textbf{profiles multiply; gates decide what may count; contractions
decide what to ask.}

\subsection{6. The open question: which contractions mean
something}\label{the-open-question-which-contractions-mean-something}

The product is not a design choice --- two profiles multiply the same
way regardless of what you intend to ask. All the design freedom lives
in the contraction: which indices to sum away, and which to leave
standing. A contracted index is a question answered in aggregate; a free
index is an axis the reader still gets to see. Every statistic in \S 4 is
one point in this space --- leave the topic index free and you get Hot
Topics; contract everything and you get a score. The general question
the sync poses, for any statistic anyone proposes adding, is simply:
which indices does it contract, and can its receipts still be shown?

One rule emerges from the cases that work: contract only the indices
whose variation you are willing to stop seeing --- and never contract an
index before a reading that needs it. (Jason's amendment, 2026-10-21:
the obvious candidate counterexample is the hypermap's own coloring ---
each category and subcategory colored by rolling its rated children up.
But categories and subcategories carry direction too: the roll-up runs
\emph{within} each direction, give-children into the parent's give,
receive into its receive, so the index actually summed is topic alone
--- children into parent --- while direction stays free and positional
at every level of the hierarchy. The heat map is therefore not an
exception to the rule but an instance of it, one level up: a contraction
over a coarser topic index, direction untouched.) Order matters,
mathematically and not just presentationally. Sum a profile's give and
receive components into a single number per topic, and you have
destroyed, before any role question is asked, exactly the components in
which roles live.

Roles are the clearest case. A role --- Dominant, Submissive, Rigger,
Brat --- is a characteristic pattern over the same (topic, direction)
components a profile carries: in the essay's language, a
\textbf{prototype vector in the same 744-dimensional space}. Reading a
person against the roles is then a contraction of one profile against
each prototype, summing over topic and direction, with the \textbf{role
index left free}. The result is a spectrum: one number per role, each
normalized by the person's total rated mass, so it reads as the fraction
of their rated self whose shape matches that prototype. (Familiar
horoscope-style role tests print exactly this spectrum with the
arithmetic hidden; here the arithmetic is the point.) And the partial
version --- topic index left free --- returns the per-topic
contributions, so every suggested role arrives with its receipts, as \S 5
demands of every number in the system.

One distinction of pairing deserves stating, because the formalism has
exactly two. The sync's Krossmatch pairs direction \emph{crossed} ---
give against receive --- because it joins two people in complementary
frames. The role reading pairs direction \emph{straight} --- the
person's give against the prototype's give --- because a prototype is
not a partner across the table but a shape in the same frame. Crossed,
for two people who might meet; straight, for one person held against a
pattern.

A discovery, recorded the day the object was named (2026-10-21), extends
the straight pairing to a second employment --- and it arrived as a
question about conversation. The Conversation readings of \S 4 are
computed over the Krossmatch: they rank what two people could talk about
\emph{as prospective partners}. But conversation has two registers.
\textbf{Crossed conversation} is negotiation talk --- what could happen
between us. \textbf{Straight conversation} is peer talk --- two riggers
comparing rope, two submissives comparing being tied --- and it lives in
exactly the components the Krossmatch discards: give against give,
receive against receive. Play cannot use that parallel sector; nothing
flows there. Talk can, because talk needs shared ground, not flow ---
and shared ground comes in both orientations. The formalism therefore
has two pairing \emph{kernels}, crossed and straight, and the straight
kernel two employments: a person against a prototype (the roles above),
and a person against a person in the peer register. The shipped
Conversation mode now reads both registers: the crossed register in
cross matches, and the straight register in straight matches, where it
is shipped as Straight Talk (\S 9). The prototype employment remains what
it always was here: a mirror, not a gate --- a reading that contracts a
single profile, speaks only to that person about themselves, is theirs
to approve or discard, and never enters the sync between people.

Two honesties attach. First, the prototypes are themselves statistics
--- hand-authored at first, refinable later from profiles their owners
donate for exactly that purpose --- and a prototype built from a thin
cohort carries real uncertainty, which an honest display encodes rather
than hides. Profiles multiply; prototypes compare; and the choice of
what to contract remains, as ever, the choice of what to ask.

\subsection{7. The heat map: a contraction up the
hierarchy}\label{the-heat-map-a-contraction-up-the-hierarchy}

The pair product is not the only contraction in the system. A single
profile contracts too, every time the hypermap colors a stem --- and
because categories and subcategories carry direction just as topics do,
the contraction runs within each direction separately, exactly as \S 6's
rule requires. No second profile takes part in it; there is no pairing
here, and no Krossmatch --- one profile, summed against its own
hierarchy.

The index summed is topic, rolled upward through the hierarchy. A
subcategory's give value is a per-direction aggregate of its topics'
give ratings, counting only topics rated in that direction --- unrated
children are excluded, not zeroed. A category's value is the same
aggregate taken over its \emph{complete} topic list: it does not average
its subcategories' averages, so every rated topic votes regardless of
how the hierarchy groups it. {[}Session: confirm whether each aggregate
is a mean or a sum --- the exact formula is unverified in this pass and
must be checked against the code before this paper states it.{]} No one
rates a subcategory or category directly --- that feature existed once
and was deprecated, and the map is better without it. Where a stored
value nevertheless sits at that level, it arrived as coarse import
residue, and the shipped code does let such a value override the
children's aggregate for display --- a survival the owner has ruled
against (see the design notes below): the per-direction aggregation is
the tapers' universal characteristic, never to be superseded by a stored
rating, and branch-level import seeds are to be resolved at topic level
through the suggested-ratings popularity list rather than by standing in
for their topics.

The aggregate is then read as a color --- continuously, not in bands.
Exactly zero is the curious yellow, pure. Positive values blend toward
green in proportion to their size, negative toward red; a stem half with
no votes at all is gray. The two directions are drawn as the two tapered
flanks of the stem itself, give on the flank above the spoke's midline,
receive below --- the same positional grammar the topic level uses. The
node discs carry no rating color at all. Their color, like the
midline's, comes from the hanky code --- the old role code, in which the
color names the interest and the pocket names the direction: left for
the one who does, right for the one who receives. The app inherits the
grammar whole: color says what, position says which way. Only the stems
carry the heat.

The shipped code departs from this clean description in three named
places, and the paper records them because each is a small failure of
\S 5's rule --- a number that cannot be unfolded the way it was made.
First, marks vote: a favorite, a limit, or a not-interested mark injects
a $\pm$100 vote into the same per-direction aggregate, without itself having
a direction, and a topic both marked and rated is counted twice --- so a
hard limit, which everywhere else in this system \emph{removes} a
component from the arithmetic, here silently drags its parent's color
red. Second, the spoke's tooltip prints a single ``heat'' figure per
topic that averages give and receive together first --- the one place in
the interface where direction is contracted away after all, answering a
coarser question than the two colors drawn beside it; that combined
figure is a known violation, ruled to be split per flank, with the give
taper carrying the give value and the receive taper the receive value.
Third, the numbers printed in tooltips, tags, and rails come from a
\emph{second} aggregator with different rules --- unrated children
dilute it, unanimity forces it to 100, and for subcategories it reads
only imported children rather than the built-in topics that color the
stem --- so the map's colors and the map's numbers are two different
statistics of the same profile, and they can disagree.

None of this invalidates the reading. The stems remain an honest
contraction: one index summed, within direction, receipts intact. But
the deviations are the same species as every defect this project has
found in itself --- a summary that quietly stopped being the summary it
claimed to be.

\emph{(Owner's design note, 2026-10-21, not yet shipped: the category
value should also consider its subcategories' averages --- each
subcategory entering its parent's computation as a unit with its own
averaged value, rather than the parent reaching past the subcategories
to a flat aggregate of topics alone. The taxonomy's groupings are real
structure; the contraction should respect them at every level, not only
the lowest.)}

\emph{(Owner's ruling on stored ratings, 2026-10-21: direct rating of
categories and subcategories is a deprecated feature and is not
returning. The per-direction aggregation over the tapers is the
universal characteristic of branch color, never superseded by a stored
value; the code's override is slated for removal, and branch-level
import seeds resolve at topic level through the suggested-ratings
popularity list.)}

\emph{(Owner's adjudication of the deviations, 2026-10-21: the map's
colors and the map's numbers must be one statistic --- two aggregators
that can disagree are a defect, not a nuance. And a tooltip belongs to a
flank, not to a spoke: the give taper carries the give value, the
receive taper the receive value. A single spoke tooltip that contracts
direction to cover both is not a meaningful object, and its combined
``heat'' figure has no question it honestly answers.)}

\subsection{8. The Krossmatch}\label{the-krossmatch}

\S 2 named the object in passing; this section examines it, because it is
the joint the whole paper stands on.

The tensor product of two profiles is everything --- every component of
one against every component of the other. The Krossmatch is what you get
by restricting that product along a fixed \textbf{pairing kernel}: the
topic index is paired with itself, and the direction index is paired
with its \emph{complement}. Formally, let $\sigma$ be the swap on directions
--- give exchanges with receive, and applying it twice returns you home:
$\sigma^2$ is the identity. The kernel is then the identity on topics tensored
with the swap on directions, and the Krossmatch's components are
M\_\{t,d\} = P\_\{t,d\} $\cdot$ Q\_\{t,$\sigma$(d)\}: for each topic, each direction,
one person's value times the other's value in the opposite direction.
Equivalently, define the \textbf{complement map} $\kappa$, which re-reads a
profile in the crossed frame --- ($\kappa$Q)\emph{\{t,d\} = Q}\{t,$\sigma$(d)\} ---
and the Krossmatch is the first profile laid component-wise against the
complement of the second. The parallel components of the product ---
give against give, receive against receive --- exist in the product and
are simply not in the object. (They are not waste: \S 6 records the peer
register of conversation that reads them.) The Krossmatch is defined by
what it refuses as much as by what it keeps.

In the language of the definitions index: the Krossmatch is the outer
product of the two profiles read through the swap kernel along the
crossed diagonal. The outer product is the product stage; the diagonal
read and its sum are the contraction stage. {[}Session: wording to
confirm.{]}

The bowtie definition, as corrected by the owner on 2026-10-26, is
stated here in his terms, with his caveat kept visible. The bowtie
\emph{is} a direct product with a kernel applied that swaps give and
receive on one profile, and the crossed contraction follows: my give
times your receive, plus my receive times your give. The straight
reading is the same direct product without the swap, with same-direction
terms. The owner states this understanding with the caveat, in his own
words, that he may be using a wrong language there. This paper therefore
records the definition as stated and flags it for the session to test,
rather than asserting it as settled mathematics if the mathematics does
not support the wording. An earlier instruction to sweep direct-product
language from the paper is withdrawn and is not applied here.

In arithmetic, with no machinery at all: a topic where Adam Peasant
stands at give 80, receive 10 and Eve Vale at give 20, receive 90
crosses to 80$\times$90 + 10$\times$20 = 7,400. Make Eve Vale a top instead --- give
90, receive 20, the same total appetite for the same topic --- and the
crossing reads 80$\times$20 + 10$\times$90 = 2,500. No rule anywhere says two tops fit
worse. The multiplication is the rule: when both people lean the same
way, every product in the sum contains a small factor, and the fit
starves on its own.

Summed whole, the Krossmatch induces a bilinear form on profiles, B(P,
Q) = $\Sigma$\_\{t,d\} P\_\{t,d\} $\cdot$ Q\_\{t,$\sigma$(d)\}, and the form has a signature
worth naming. Rotate the direction plane by 45 degrees, into the
\textbf{blend} axis --- (give + receive)/$\sqrt{2}$ --- and the
\textbf{difference} axis --- (give - receive)/$\sqrt{2}$ --- and the form reads
B = blend$\cdot$blend - difference$\cdot$difference. In direction space, per topic,
its signature is (1, -1): indefinite, the way spacetime's interval is
indefinite, though nothing here is a spacetime. The decomposition says
in one line what the arithmetic example says in numbers: agreement in
overall appetite adds to a pair's fit; agreement in directional lean
subtracts from it. Opposites fit --- not as a slogan, but as the
signature of the only pairing the sync performs.

The summed form is symmetric: exchanging the partners leaves B
unchanged, because the sum runs over every topic and both directions.
The flanks are not symmetric --- exchanging the partners swaps the
give-to-receive channel with the receive-to-give channel --- and \S 10's
dark sums are exactly the instruments that hear that asymmetry. Symmetry
of the sum, asymmetry of the channels: both facts, priced separately.

Physicists will recognize the shape of the choice. Indices can be
contracted through an identity, which pairs each thing with itself, or
through a pairing tensor, which pairs each thing with its dual ---
spinor indices are contracted through the antisymmetric $\epsilon$, not through a
metric, and that choice is what makes a spinor's partner its complement
rather than its copy. The Krossmatch is the same species of decision,
made once and made everywhere: directions are contracted through the
swap, so that each direction's dual is the other. Give has no meaning in
this algebra except against receive, and the pairing is where that is
enforced.

Two separations keep the object honest. First, the Krossmatch is
\textbf{not gated}: it is pure bilinear algebra, the raw grid of flows.
The weaker-side rule, the route table, the limits, Heat --- \S 3's whole
apparatus --- are nonlinear semantics applied to its components
afterward, and conflating the two layers would make the algebra
impossible to audit. Second, the Krossmatch is \textbf{one of exactly
two pairing kernels} in the formalism. The other kernel pairs direction
straight, and it carries two employments: the role reading of \S 6, where
a prototype is a shape in the same frame rather than a partner in the
opposite one, and the peer register of conversation, where the product's
parallel sector is read as shared ground between two people. Everything
else in the system that looks like a pairing --- the heat map of \S 7 ---
is no pairing at all, but a single profile summed against itself.

A scope note, for the same reason: as constructed, the Krossmatch is
pairwise. Group readings (\S 4) gather pair Krossmatches under majority
rule, and the construction plainly suggests the generalization --- each
profile crossed against the rest --- but a single group-native object,
one tensor for \emph{n} profiles, is not claimed here as well-defined.
The pair object is exact; the group form works theoretically, and the
paper keeps the distinction visible.

That is the object's power and its limit. It says nothing about what a
pair should do; it is the complete, un-gated statement of what could
flow between them, computed by the one arithmetic that refuses,
structurally, to pretend that wanting the same side is the same as
fitting. Every verdict the sync delivers is a gated contraction of this
grid --- and when a verdict is questioned, this is the grid it answers
to.

\subsection{9. The Straight Match}\label{the-straight-match}

\S 8 closed by separating the Krossmatch from the only other pairing
kernel in the formalism. This section gives that kernel its object,
because it has one --- as definite as the Krossmatch, and disjoint from
it.

Take the product of two profiles and look at one topic at a time. Where
a topic meets itself, the product's components form a 2$\times$2 block over the
direction pairs: give against give, give against receive, receive
against give, receive against receive. The block has exactly two
sectors. The \textbf{diagonal} --- give against give, receive against
receive --- pairs each direction with itself: the \emph{parallel}
sector. The \textbf{anti-diagonal} --- give against receive, receive
against give --- pairs each direction with its complement: the
\emph{crossed} sector. The two sectors share no component, and between
them they exhaust the block: per topic, the product's self-pairing is
exactly these two sectors and nothing else.

The Krossmatch is the anti-diagonal, kept by the pairing kernel of \S 8:
identity on topics, swap on directions. The diagonal is kept by the
mirror-image choice --- identity on topics \emph{and} identity on
directions, pairing every index with itself through the Kronecker delta,
where the Krossmatch pairs direction through the swap. The object it
defines is the \textbf{Straight Match}, written with the circled dot:
Adam Peasant $\odot$ Eve Vale. Its components are S\_\{t,d\} = P\_\{t,d\} $\cdot$
Q\_\{t,d\} --- for each topic, each direction, one person's value times
the other's value \emph{in the same direction}. Algebraists know the
operation as the Hadamard product: the component-wise product of two
arrays read in the same frame. Where the Krossmatch lays a profile
against the complement of another, the Straight Match lays it against
the other exactly as it stands.

The two objects are one operation viewed twice, and the notation can say
so: \textbf{P $\bowtie$ Q = P $\odot$ $\kappa$Q} --- the Krossmatch is the Straight Match
computed with the second profile crossed first, by the complement map of
\S 8. There is a single primitive here --- multiply matching components
--- and a single frame choice, straight or crossed, that decides what
\emph{matching} means. A pairing is a choice of what counts as a match;
everything else in this section is bookkeeping for that choice. In the
language of the definitions index: the straight reading is the
same-direction diagonal, inner-product-like; the crossed reading is the
outer product through the swap kernel. {[}Session: wording to
confirm.{]}

Sum each object whole and the bookkeeping pays out in two different
currencies. The Straight Match contracts to the ordinary inner product
$\langle$P, Q$\rangle$ --- a positive-definite form, the Euclidean geometry of the
profile space itself. It is a \emph{similarity} measure, and agreement
feeds it without exception: two components pointing the same way
multiply to a positive term whether what is shared is appetite or
aversion, because in a single frame likeness is likeness --- in lean as
much as in appetite. The Krossmatch contracts to B(P, Q) of \S 8,
indefinite, signature (1, -1): agreement in appetite adds, and agreement
in lean subtracts. Similarity and fit are different facts, and the
formalism now prices them separately. Two tops carrying \emph{identical}
maps are the limiting case: maximally similar --- the Straight Match
reads as high as the scale allows --- and minimally fit, their crossing
starving on small factors exactly as \S 8's arithmetic showed. A single
compatibility number would have to call that pair both perfect and
hopeless. The two pairings decline to confuse the two statements; there
is a number for each fact.

The Straight Match is no longer idle. \textbf{Straight Talk} is its
shipped reading between two people, in Conversation mode only and in
straight matches only, never in Play and never in a crossed match: for
each shared topic, the straight sum --- give$\times$give plus receive$\times$receive,
limits excluded, a null sum where neither direction is known on both
sides --- ranked highest first, ties by topic name, and banded over
equal thirds of the per-topic range (-20,000 to +20,000): Debate at
-6,667 and below, On-the-fence from -6,666 to +6,666, Shop talk at
+6,667 and above, no band for a null sum. The rings wear the band names
with their counts, and a Limits section stands beside them, display
only. What Straight Talk reads is what two people hold in common, in the
same orientation --- the conversation register \S 6 discovered in the
sector the sync had always discarded, now computed in the open.
Peer-ground is dropped as a term; it is not used in this revision. The
prototype employment stays where \S 6 put it: a mirror reading against a
pattern, not a gate, and not a shipped score.

Straight Talk's headline follows \S 4. Debate topics count in the Parlay
MS numerator and denominator, like Shop talk and On-the-fence topics.
Bands display sections only; none of the four pair numbers refers to
banded topics. {[}Flag: provisional, as in \S 4.{]}

Straight Talk's list is banded, not merely ranked. Debate topics join
the list. The list splits into band sections, Shop talk first, then
On-the-fence, then Debate, covering the full band ranges, with topics
ordered within each section by straight sum descending. Limit-blocked
topics stay out of the sections; their home is the Limits section.
One-sided topics, with no sum, stay out of the sections as well. This
paper states the mechanism and the bands; it cites no straight census,
so no census claim in this section depends on resolving the one-topic
difference between a paper computation and a live readout for a
73-shared-topic pair, which remains unresolved in the companion notes.

The straight rings follow \S 4. On a circledot tab, the inner ring is the
Parlay Kink Factor and the outer ring is the Parlay Match Score, both
provisional in definition. The Parlay mirror is carried in \S 4 with its
flags; this section does not harden it by repeating it without those
flags.

\subsection{10. The Penumbra}\label{the-penumbra}

\S \S 8 and 9 leave no third sector hiding in the product. For each topic,
the product's 2$\times$2 direction block splits into the diagonal and
anti-diagonal sectors --- parallel and crossed, disjoint and together
exhaustive. The penumbra begins one level up, when those components are
summed. A sum can read a sector with a plus sign, or read it against
itself; both choices are arithmetically real, and only one of them is
lit.

Suppress the topic index and write the four products as \emph{gg},
\emph{rr}, \emph{gr}, and \emph{rg}: give against give, receive against
receive, give against receive, receive against give. There are four
independent signed sums. Two are already the paper's working quantities:
\textbf{S = gg + rr}, the straight sum, whose contraction is
resemblance; and \textbf{C = gr + rg}, the crossed sum, whose
contraction is fit. The engine reads this lit pair. The other two are
\textbf{K = gg - rr}, the channel imbalance of a resemblance, and
\textbf{$\epsilon$ = gr - rg}, the orientation of action-flow. K is unchanged
when the partners are exchanged. $\epsilon$ is not: exchanging them exchanges
\emph{gr} and \emph{rg}, so $\epsilon$ changes sign. It records which way the
current runs, a fact no symmetric sum can contain.

Darkness, then, is not a property of hidden components. The four
products sit in the product in plain sight. It is a choice of signs at
the moment of summation --- which combinations an instrument reads, and
which it leaves standing unread in the penumbra.

The receipt is arithmetic. Let one person stand at give 80, receive 20,
and the other at give 30, receive 90. The crossed products are 80 $\times$ 90 =
7,200 and 30 $\times$ 20 = 600, so C = 7,800 and $\epsilon$ = 6,600. Reverse the
directional lean --- give 30, receive 90 against give 80, receive 20 ---
and the same two products change places: C is still 7,800, while $\epsilon$ is
-6,600. Addition cannot tell the pairs apart; the dark sum can. Same
crossed fit, opposite action-flow.

K answers a different question. Two people alike chiefly as givers, each
at give 90 and receive 10, have S = 8,200 and K = 8,000. Two people
alike chiefly as receivers, each at give 10 and receive 90, have the
same S and K = -8,000. The straight sum rightly calls the resemblances
equal. K says that they live in opposite channels --- a distinction peer
talk would hear immediately, even though no shipped score reports it.

\textbf{Heat as a phase.} The lit sums are timeless: they describe two
profiles, not the occasion on which the profiles meet. Heat was salvaged
as a pairing-level weather by scoping it carefully. It does not
complexify a profile. Saved profiles remain lists of real magnitudes,
exactly as \S 1 defines them. At a pairing, each person wears one
\emph{declared} angle for the session --- a garment over the profile,
not a change to it --- and the pairing is read in the relative angle.
Under the sesquilinear reading, the lit scores pick up the alignment
factor cos($\Delta$$\theta$), where $\Delta$$\theta$ is the difference of the two angles. A single
angle is unobservable alone; only the difference enters any score, which
is why Heat belongs to the pairing and cannot be stored as one person's
trait.

The clock is global: one estrus clock per person, with two scales. The
slow hand is the estrous cycle, turning over weeks. It may be informed
by a tracker, and it modulates the fast hand --- Heat proper, the
declared angle of tonight, turning over days. Where a tracker and a
declaration disagree, the declaration takes precedence: the person is
the authority on their own state, and a tracker is an informed prior,
never an imposed reading. At the Clarity seam, the fast hand does not
glide through zero; desire resets by a jump, sometimes into aversion.
When one partner has crossed that seam and the other has not, $\Delta$$\theta$
approaches $\pi$, cos($\Delta$$\theta$) approaches -1, and the pairing \emph{inverts}
rather than merely cancelling. The shipped Heat slider is the fast
hand's ancestor: a session-only, declared threshold state, never written
into a saved profile, reporting the same weather in a coarser
instrument.

One hand remains undecided, and the section leaves it open on purpose.
Arousal --- the within-the-evening variable --- might be a third phase:
then the session garment is a 3-torus, arousal carries an angle of its
own, and it could invert an alignment by itself. Or arousal might be a
gain on the heat hand: then the garment is a 2-torus with an amplitude,
arousal can raise or lower the size of a reading but never reverse its
sign. The mathematics permits both; which one the body is doing is an
open question, and the formalism will not pretend otherwise.

\textbf{Initiative.} Not every named quantity in the system is a
contraction, and Initiative needs no second algebra. It is bookkeeping
at the resolution layer --- computed only after \S 3's gates have decided
which channels are live, and attributed per channel rather than summed
away. For a live channel from a giver to a receiver, define the
\textbf{surplus} as the giver's give minus the receiver's receive.
Rectify it: a positive surplus attributes initiative to the giver, whose
wanting overflows the appetite it meets; a negative surplus attributes
initiative of the same magnitude to the receiver, whose appetite summons
the scene. The totals keep their territory partition --- a surplus
earned on one topic is not cancelled by an opposite surplus on another,
the way $\epsilon$'s terms cancel. Initiative says \emph{where} each person's
wanting is doing the driving, topic by topic.

The worked check is the one that forced the distinction. In a
teacher/student scene, on the discipline channel, let the teacher's give
be 70 and the student's receive be 95. The surplus is 70 - 95 = -25:
initiative to the student, 25. The teacher may hold the authority of the
role entire; the scene's unfolding can still be driven by the student's
appetite. Authority and initiative are different facts, and this
bookkeeping is the first instrument in the design that can say so in
numbers.

Initiative is therefore named and formalized here as a \emph{candidate},
not built in the app. Two questions remain the owner's to settle before
it can be more: whether surplus is the right attribution rule at all,
against the other rules the record sketches; and whether initiative
earned on Curious-fit channels --- the explorative frontier, where the
wanting is newest --- should carry extra weight.

\subsection{11. The multi-pole
extension}\label{the-multi-pole-extension}

\S \S 1--10 are written for the founding domain's topics, each of which
carries exactly two components --- give and receive --- and the two
pairing kernels of \S \S 8--9 are exact there. The taxonomy has since grown
a second species of topic, and this section records how far the
formalism honestly extends to it, including where it presently does not.

A \textbf{pole topic} carries from one to eight poles --- stances, not
directions. A triangle topic (three poles) is rated by placing a point
in the simplex of the three stances; a radar topic (five to eight poles)
is rated pole by pole, each an independent affinity; a diamond carries
four. A person's rating of a pole topic is therefore a vector of
per-pole affinities, each from -100 to +100, and a zero is a rating, not
an absence. The engine compares two such vectors, as shipped, by
\textbf{agreement}: 100 minus half the mean absolute difference over the
poles --- a number that reads as a percentage of sameness of shape. Each
side also carries a \textbf{Depth Factor} intensity for the topic ---
the vanilla skins' analogue of the Kink Factor --- defined as the
largest absolute pole affinity that side holds: how deep the topic goes
for that person, set by its strongest stance. Pole ratings count as
ratings everywhere a two-flank rating does, and the straight register's
considered set includes them.

Three findings from the code trace of 2026-11-19 are recorded here
because each is a place where the extension falls short of this paper's
own standard, and the paper's standard is that shortfalls are stated,
not smoothed.

First, \textbf{the two kernels have not yet been given distinct
constructions for pole topics.} A pole topic stores one vector per
person; there is no complement map on stances analogous to \S 8's swap on
directions, crossed fits for pole topics are null by construction, and
the straight reading at present substitutes the crossed numbers
identically. For two-flank topics the crossed and straight objects are
disjoint sectors of one product, priced by forms of opposite signature;
for pole topics the two registers currently \emph{coincide}, and the
distinction the formalism exists to protect is, for this species,
unimplemented rather than false. The owner's tentative semantics stand
recorded as tentative: the crossed register as \textbf{communion} ---
Play in this domain; living, working, and being together in the others
--- and the straight register as \textbf{an interesting conversation}. A
construction honouring that distinction for poles is an open design
problem, and one proposed fork is recorded in \S 12.

Second, \textbf{the Depth Factor's intensity is sign-blind.} Being
computed from absolute affinities, it rises identically for a shared
enthusiasm and a shared aversion: two people who both stand at -90 on a
stance register the same depth as two who stand at +90. \S 9's straight
form has the same property and there it is \emph{correct} --- in a
single frame, likeness is likeness, in aversion as in appetite. In an
intensity that purports to say how deep a topic goes for a pair, the
discarded sign is a fact the reading should not lose, and the paper
flags the Depth Factor's present form as a defect against \S 5's rule: a
number whose receipts, opened, do not show the thing the number was
taken to mean.

Third, \textbf{pole topics do not yet appear in either register's match
or curious lists}, though they enter headline aggregates and the
Scribe's pair readings; and in group readings of three or more, their
depth entries at present compute to zero. These are implementation gaps,
recorded as such. A reader of the shipped numbers should know that the
pole species is counted in the totals and absent from the lists.

\subsection{12. Intensity as an independent
coordinate}\label{intensity-as-an-independent-coordinate}

The ratings of \S 11 carry one further component, added by the owner's
ruling of 2026-11-19, which this section states as a model and then,
separately, as an implementation status --- because the two are not the
same, and the paper does not pretend otherwise.

\textbf{The model.} For a topic of \emph{n} poles, the leanings are
\emph{n} numbers: the blend, where the person stands among the stances.
\textbf{Intensity is an (n+1)th number}, stored alongside the leanings
and not derived from them: it is the topic's \textbf{relevance} --- how
much the person cares about the issue at all --- and it is signed.
Positive intensity is caring in the issue's favour; zero is
indifference, the don't-care point; negative intensity is \emph{aversion
to the issue itself}, a stance toward the question rather than within
it. The display rule was ruled the same day: intensity shows as
\textbf{brightness} --- full brightness at maximum positive, dimming to
don't-care at zero, and passing into complementary colours below zero.
The ruling supersedes the earlier reading, in which intensity was
computed from the leanings (as the Depth Factor of \S 11 still computes
it): two people can hold identical blends with utterly different
relevance, and a formalism that derives the second from the first has
destroyed, before any question is asked, exactly the component in which
caring lives --- the same species of error as contracting direction in
\S 6's rule, one index over.

\textbf{The status.} The implementation is partial, and its partiality
is specific. The pass shipped on 2026-11-19 decoupled the
\emph{drawing}: a pole widget's shape now holds still while its
intensity is swept, where before the shape was the intensity's prisoner.
But storage still derives intensity from the leanings, and the intensity
control still rescales stored values --- the halfway house the ruling
was issued to replace. The full build is specified: intensity stored as
its own number per topic, existing profiles seeded from the current
derivation so no one's map changes meaning under them, the intensity
control editing that number alone. It stands first in the current
development handover. One question is expressly reserved: \textbf{how
the (n+1)th coordinate weights pair scoring} --- whether relevance
multiplies agreement, gates it, or stands beside it as a third reported
number --- is to be settled by a proposal put to the owner, not assumed
by the implementers, and this paper records it as open.

\textbf{A fork, recorded.} \S 11's open problem --- a crossed construction
for poles --- has one proposed route that touches this coordinate. If
each pole's sign were freed to carry direction (positive for giving,
negative for receiving, in the founding domain's grammar), a swap-like
kernel on poles would exist and the crossed register could be built. The
owner's own objection is recorded with the proposal, because it may be
fatal: one sign cannot carry both aversion and direction at once. The
fork is therefore: either aversion lives \emph{only} on the (n+1)th
coordinate's sign --- freeing the pole signs for direction --- or the
pole signs keep their present meaning and the crossed construction must
be found elsewhere. Unruled, and stated here so that whichever way it
settles, the losing consideration is on the record.

\textbf{Not Heat.} The (n+1)th coordinate is a stored, standing property
of a person's rating of a topic. It is not \S 10's Heat, which is a
declared, per-session state of a \emph{pairing}, never written into a
profile. Relevance is what a person brings; Heat is the weather of the
occasion on which two profiles meet. The system keeps both, and this
paper keeps them distinct.

\emph{End of the 2026-11-21 revision: \S \S 11--12 as added above. \S \S 1--10
and Appendix A stand as revised 2026-11-06. Open items are marked
{[}Flag{]}, {[}Session{]}, or as recorded findings in place.}

\emph{End of revision working draft, 2026-10-29. Front matter,
definitions index, and \S \S 1--10 as revised above. Open items are marked
{[}Flag{]} or {[}Session{]} in place and are listed with status in the
companion notes for 2026-10-29. No headline in this draft is a
kinkiness-weighted recomputation, no Personal Kink Factor number is
asserted, and no group weighted-Kink Factor example is invented.}

\subsection{Appendix A. The six numbers, side by side ---
HELD}\label{appendix-a.-the-six-numbers-side-by-side-held}

\textbf{HELD appendix, not final.} This appendix is held for the full
review sometime after Saturday. It is not the interim table for the
munch, and inclusion here does not harden anything it contains. Items
marked interim, provisional, or recommendation below stay interim,
provisional, or recommendation. Where this appendix corrects the held
two-table document it was copied from, the correction is stated in the
notes at the end of this appendix, and the paper's ruled values govern.

Source: the held two-table document delivered 2026-10-29, Table 1 (Kink
Factors) and Table 2 (outer numbers), with its page-3 key and open list,
as corrected by the 2026-10-30 provisional six-score rulings in \S 4.
Values are copied from that document, except for the corrections in the
notes below. The six columns are the six numbers of \S 4: Personal and
Charted for one person, Play KF and Play Match Score for a crossed pair,
Parlay KF and Parlay Match Score for a straight pair.

\subsubsection{Table A1. Kink Factors (inner
rings)}\label{table-a1.-kink-factors-inner-rings}

{\def\LTcaptype{none} % do not increment counter
\begin{longtable}[]{@{}
  >{\raggedright\arraybackslash}p{(\linewidth - 6\tabcolsep) * \real{0.2500}}
  >{\raggedright\arraybackslash}p{(\linewidth - 6\tabcolsep) * \real{0.2500}}
  >{\raggedright\arraybackslash}p{(\linewidth - 6\tabcolsep) * \real{0.2500}}
  >{\raggedright\arraybackslash}p{(\linewidth - 6\tabcolsep) * \real{0.2500}}@{}}
\toprule\noalign{}
\begin{minipage}[b]{\linewidth}\raggedright
Row
\end{minipage} & \begin{minipage}[b]{\linewidth}\raggedright
Personal Kink Factor
\end{minipage} & \begin{minipage}[b]{\linewidth}\raggedright
Play Kink Factor
\end{minipage} & \begin{minipage}[b]{\linewidth}\raggedright
Parlay Kink Factor
\end{minipage} \\
\midrule\noalign{}
\endhead
\bottomrule\noalign{}
\endlastfoot
Tab / pair & Profile tab, individual & Bowtie, Play pair & Circledot,
Conversation pair (straight match) \\
Status & Provisional (2026-10-30), with open details & Provisional
(2026-10-30) & Provisional (2026-10-30) \\
Register & Whole map, individual & Play & Conversation only, not Play \\
Counted & Per topic: maximum flank (signed, proportional,
epsilon-thresholded) plus flag sum once, weighted. Not a counted /
considered ratio --- an additive expanse. Not clamped at 100. &
Geometric mean over members of (binary rated-or-not + table sum), times
maximum Krossmatch fit, normalized by group maximum & Same shape on the
Straight Match, with the Parlay table and considered set \\
Considered & All 372 topics. Limit-only topics count nothing; the
denominator stays 372. & Every shared, non-limit-blocked topic, weighted
& Every shared, non-limit-blocked topic, weighted. Fantasy included. \\
Rating basis & Maximum flank $\times$ hierarchical share $\times$ kinkiness, per
topic. A negative maximum can pull a topic negative. & Cross: give $\times$
receive, normed /100. Binary fit iff \textbar term\textbar{}
\textgreater{} epsilon. & Straight sums: give $\times$ give + receive $\times$
receive, normed /100. Binary fit iff \textbar term\textbar{}
\textgreater{} epsilon. Bands at $\pm$6,667 are display only (see Key
below). \\
Limits / one-sided & Limit-only topics count nothing in both rings;
denominator stays 372 as the penalty for a hard limit. Limit plus
another mark is excluded altogether. No -1 limit entry in the table. &
Limit-blocked topics out of considered & Limit-blocked out; one-sided
topics out of counted; a null sum stays in considered, not counted \\
Weights / groups & Hierarchical shares $\times$ kinkiness, once. & Hierarchical
shares $\times$ kinkiness, pairs; strict majority counting, groups, normalized
by group maximum & Hierarchical shares $\times$ kinkiness, pairs; strict
majority counting, groups, normalized by group maximum \\
\end{longtable}
}

\subsubsection{Table A2. Outer numbers: Charted and the Match
Scores}\label{table-a2.-outer-numbers-charted-and-the-match-scores}

Each column here is the outer partner of the same-colour column in Table
A1: Personal with Charted on the profile tab, Play Kink Factor with Play
Match Score on the bowtie, Parlay Kink Factor with Parlay Match Score on
the circledot.

{\def\LTcaptype{none} % do not increment counter
\begin{longtable}[]{@{}
  >{\raggedright\arraybackslash}p{(\linewidth - 6\tabcolsep) * \real{0.2500}}
  >{\raggedright\arraybackslash}p{(\linewidth - 6\tabcolsep) * \real{0.2500}}
  >{\raggedright\arraybackslash}p{(\linewidth - 6\tabcolsep) * \real{0.2500}}
  >{\raggedright\arraybackslash}p{(\linewidth - 6\tabcolsep) * \real{0.2500}}@{}}
\toprule\noalign{}
\begin{minipage}[b]{\linewidth}\raggedright
Row
\end{minipage} & \begin{minipage}[b]{\linewidth}\raggedright
Charted
\end{minipage} & \begin{minipage}[b]{\linewidth}\raggedright
Play Match Score
\end{minipage} & \begin{minipage}[b]{\linewidth}\raggedright
Parlay Match Score
\end{minipage} \\
\midrule\noalign{}
\endhead
\bottomrule\noalign{}
\endlastfoot
Tab / pair & Profile tab, individual, outer ring & Bowtie, Play pair,
outer & Circledot, Conversation pair, outer \\
Status & Provisional (2026-10-30), with open details & Provisional
(2026-10-30) & Name ruled; definition provisional (2026-10-30) \\
Register & Whole map coverage, individual & Play & Conversation only,
not Play \\
Counted & Attended topics: binary rated (epsilon-thresholded, either
flank) plus flag sum once, weighted. Charted counts topics, not flanks.
& Geometric mean over members of normalized (binary fit + table sum) &
Same shape on the Straight Match, with the Parlay table and considered
set. Debate topics count in numerator and denominator. \\
Considered & All 372 topics. Limit-only topics count nothing; the
denominator stays 372. & Play-considered: considered minus fantasy-scope
topics (fantasy is out of both numerator and denominator) & Every
shared, non-limit-blocked topic. No fantasy subtraction --- Fantasy is
included in the Parlay considered set. \\
Rating basis & Binary, not a percentage. A topic counts when rated
(epsilon-thresholded) or flagged; negatives count by absolute value to
binary. Flag weights sit on top (see Shared contributions below). &
Cross: give $\times$ receive, normed /100. Binary fit iff
\textbar term\textbar{} \textgreater{} epsilon. & Straight sums, as in
Parlay Kink Factor, normed /100. \\
Limits / one-sided & Limit-only topics count nothing in both rings;
denominator stays 372. & Limit-blocked topics out of considered &
Limit-blocked out; one-sided topics have no sum and no band \\
Weights / groups & Hierarchical shares $\times$ kinkiness. Topics, not flanks.
& Hierarchical shares $\times$ kinkiness, pairs; strict majority counting,
groups, normalized by group maximum & Hierarchical shares $\times$ kinkiness,
pairs; strict majority counting, groups, normalized by group maximum \\
\end{longtable}
}

\subsubsection{Shared contributions --- Personal and
Charted}\label{shared-contributions-personal-and-charted}

Same table, per topic, once, times hierarchical share $\times$ kinkiness.
Multiple flags on one topic all add, once.

\begin{itemize}
\tightlist
\item
  Tried: +1 for the entry.
\item
  Fantasy: +0.5.
\item
  Teased: 0. Curious: 0. Both were removed from this table.
\item
  Favorite: weight 1, counts with an additive flag.
\item
  Into it: adds nothing on its own (it is the rating condition, not an
  additive flag).
\item
  Limit: the -1 entry was removed from the table altogether. Limit-only
  topics count nothing instead (both rings, denominator 372).
\end{itemize}

In Personal the rating itself is the maximum flank, signed and
proportional. In Charted the rating counts binary, epsilon-thresholded,
with these flag weights on top, as a straight sum of the table entries.

\subsubsection{Parlay table}\label{parlay-table}

Positive entries only, no null or zero entries, no banded topics.
Favorite 1, Tried 1, Fantasy 1 (full 1). Teased does not enter ---
absent, not zero. The Play table is, for now, copied from the Charted
table above; its final form is deferred. {[}Flag: provisional, final
inclusion tables deferred.{]}

\subsubsection{Key --- banded topics}\label{key-banded-topics}

A banded topic is a shared topic with a straight sum: give $\times$ give +
receive $\times$ receive, ranging from -20,000 to +20,000. Do not take a max.

\begin{itemize}
\tightlist
\item
  Shop talk: at or above +6,667.
\item
  On-the-fence: -6,666 to +6,666.
\item
  Debate: at or below -6,667.
\item
  A null sum has no band.
\item
  Limit-blocked topics go to Limits, not to a band.
\item
  One-sided topics have no sum and no band.
\end{itemize}

\subsubsection{Similarities --- and the differences that
stay}\label{similarities-and-the-differences-that-stay}

Same in both tables: the row set, the limit exclusion, hierarchical
shares $\times$ kinkiness, strict majority counting for groups in every pair
column, normalized by the group maximum, the straight bands shared by
Parlay Kink Factor and Parlay Match Score as display only, and the tab
and considered base inside each inner / outer pair.

Different, on purpose: individual versus pair shared topics; flanks
(Personal) versus topics (Charted); an additive expanse versus a
counted-over-considered ratio; binary rating (Charted) versus percentage
rating (Personal).

Inside each inner / outer pair the outer ring covers more territory,
which is why it sits above the inner, not because of any clamping.

\subsubsection{Appendix notes --- corrections to the source document,
and what stays
open}\label{appendix-notes-corrections-to-the-source-document-and-what-stays-open}

\begin{enumerate}
\def\labelenumi{\arabic{enumi}.}
\tightlist
\item
  \textbf{O6 source, corrected.} The source document says the kinkiness
  source is still open, with effective weight applying ``once the
  kinkiness source is chosen''. That wording is stale and is NOT carried
  over. The source list is RULED, seven sources: Joyal et al.~2015,
  Lehmiller 2018, FetLife, Reddit, community ratings, own survey, and
  literature priors. The strength model is RULED: anchors at factor 1.0,
  discounted sources at 0.5, log weight ln(size) $\times$ factor, prior weight
  0.5 $\times$ ln(5) $\approx$ 0.80, ignore-missing aggregation, Type 1 submissions
  excluded. What remains open is Reddit counts (pending) and the
  provisional literature priors. Whether the shipped Popularity List
  sort key implements the aggregate is unverified.
\item
  \textbf{Peer-ground, dropped.} Peer-ground was carried in the
  2026-10-29 draft as an open question. It is DROPPED in the 2026-10-30
  provisional definitions and is not a term in this revision. Bands
  display sections only; none of the four pair numbers refers to banded
  topics.
\item
  \textbf{O5 items, as provisionally ruled 2026-10-30.} Across flanks,
  the maximum flank governs Personal, not a sum; Charted counts a topic
  once. Unrated plus marked: base 0, flag still counting, limit-only
  exclusion overriding; a limit plus another mark is excluded
  altogether. Red flank with a flag: the rating fraction carries the
  sign, flag additions stay positive, net can go negative. No clamp on
  Personal or Charted. Charted per-topic flag sum counts topic flags
  once per topic, rated binary if either flank is rated, subject to an
  editor check of flank-level against topic-level behaviour in code,
  which this paper does not assert. With no recommendation yet: group
  Charted, where the record is individual only.
\item
  \textbf{Names and status, kept as flagged.} Parlay Match Score is the
  ruled name for the straight outer number, superseding Straight Match
  Score. Play Kink Factor supersedes Party Kink Factor, which is
  retired. All six definitions are provisional, 2026-10-30, for
  consideration, until the review after Saturday; final inclusion tables
  are deferred. The source document's read-across line claiming Party
  and Parlay ``match on Register family'' is not carried over verbatim,
  because Play is Play and Parlay is Conversation only --- the registers
  differ, and the similarity that does hold is inside each inner / outer
  pair, as stated above.
\item
  \textbf{No hardening by inclusion.} Nothing in this appendix changes
  the status of any item in \S 4 or in the companion notes. The Oct 28
  revision and draft 4 remain untouched. The Samhain review text is this
  same Oct 29 draft by path, so this appendix will be in Saturday's
  review text without any separate repoint.
\end{enumerate}

\emph{End of Appendix A, held for the full review sometime after
Saturday.}

\end{document}


\chapter{The Defense: Anticipated Objections and Responses}
\emph{Prepared for the candidate and the committee, 2026-11-21. Each
objection is stated at its strongest --- a defense rehearsed against
weak objections is no defense at all. Responses distinguish three kinds
of answer: \textbf{rebuttal} (the objection fails), \textbf{concession
with boundary} (the objection lands on something the work does not
claim), and \textbf{open item} (the objection identifies real unfinished
work, which is named and dated rather than disguised).}

\begin{center}\rule{0.5\linewidth}{0.5pt}\end{center}

\section{Objection 1 --- ``This is software, not
scholarship.''}\label{objection-1-this-is-software-not-scholarship.}

\emph{The candidate has built a web application. However ingenious, a
product is not a contribution to knowledge.}

\textbf{Rebuttal.} The dissertation does not ask the committee to
examine a product; it asks it to examine a \emph{framework} --- a
measured stance-space, an intensity coordinate independent of leaning,
two formally distinct comparison registers, and a receipts discipline
for aggregate claims about persons --- which happens to exist in
executable form. Research through design is an established mode
precisely for questions of the form ``can X be given workable form?'',
where the artefact is the existence proof and the dissertation is its
extraction into claims. Chapter One states the claims against the
literature; the companion paper states the mathematics. If the framework
were \emph{only} software, it could not be wrong in interesting ways ---
and it demonstrably can be: the work's own record includes a scoring
projection that failed its calibration fixtures, a register collapse
found by trace, and two aggregation defects, each identified \emph{as} a
failure because the framework made the failure stateable. Products do
not falsify themselves; research instruments do.

\section{Objection 2 --- ``\,`Tensor' is borrowed prestige. The
mathematics is
metaphorical.''}\label{objection-2-tensor-is-borrowed-prestige.-the-mathematics-is-metaphorical.}

\emph{A profile is a list of numbers; calling its pairwise products a
tensor product dresses bookkeeping in physics vocabulary.}

\textbf{Concession with boundary, then rebuttal.} The companion paper
concedes the boundary itself, in its own text (\S 5): there is no
manifold, no metric tensor, no covariance; topics are not coordinates.
What remains after the concession is not metaphor but algebra, and the
algebra does work that bookkeeping cannot. The product stage constructs
the full grid before any statistic is chosen, which is what makes every
later number auditable to its components. The two pairing kernels are
genuinely distinct operations with genuinely distinct contracted forms:
the crossed form is \emph{indefinite} --- signature (1, -1) in the
direction plane --- so that agreement in appetite adds to fit while
agreement in directional lean subtracts from it; the straight form is
the positive-definite inner product, where likeness is likeness
regardless. ``Two tops are maximally similar and minimally fit'' is a
theorem of this construction, not a slogan, and no single compatibility
scalar can state it. The vocabulary is earned where it is used and
disclaimed where it is not --- which is the correct scholarly posture
toward analogy, and the paper models it.

\section{Objection 3 --- ``Your six scores are, by your own account,
provisional. You are defending unfinished
work.''}\label{objection-3-your-six-scores-are-by-your-own-account-provisional.-you-are-defending-unfinished-work.}

\textbf{Concession with boundary.} The six score definitions stand as
provisionally ruled (2026-10-30) and the dissertation says so wherever
they appear. But the defended claims do not include the final form of
those six numbers. The claims are the stance-space, the (n+1) intensity
model, the crossed/straight distinction, the gating architecture, and
the receipts discipline --- all of which are finished enough to be
\emph{wrong}, which is the relevant standard. A doctorate defends a
contribution, not a product line's completion; and a candidate who
hardens provisional definitions for the sake of the defense would be
committing the exact offence --- a summary claiming to be what it is not
--- that the work exists to refuse.

\section{Objection 4 --- ``The taxonomies are one man's opinions with
extra
steps.''}\label{objection-4-the-taxonomies-are-one-mans-opinions-with-extra-steps.}

\emph{Sixteen domains, a thousand topics, every pole articulation chosen
by the investigator. The instrument measures the world as Jason Wrinkle
carves it.}

\textbf{Concession with boundary, then rebuttal.} The selection of
topics does encode one investigator's judgement; Chapter Two states this
as a limitation (\S 2.8) rather than discovering it under questioning.
Three boundaries apply. First, the \emph{statement} of positions is
disciplined --- the hovertip standard (each pole stated as its holders
would state it) and the neutral umbrella rule (no topic name may smuggle
a verdict) are auditable rules, and the audit record exists. Second, the
engine/skin separation makes the taxonomy \emph{replaceable}: any
scholar who believes a domain is carved wrongly can author a rival skin,
and the mathematics will process it identically --- the framework's
claims survive the replacement of every taxonomy it ships with, which is
the operational meaning of the distinction between an instrument and an
opinion. Third, single-authorship is also the source of the corpus's
unusual coherence --- sixteen domains under one fairness discipline is a
controlled comparison that a committee-authored taxonomy could not have
produced. The candidate concedes the bias and claims the control.

\section{Objection 5 --- ``There is no validation. One paying member;
262 visitors; no evidence anyone's argument ever
improved.''}\label{objection-5-there-is-no-validation.-one-paying-member-262-visitors-no-evidence-anyones-argument-ever-improved.}

\textbf{Concession, plainly.} This is the strongest objection available
to the committee, and the candidate does not dilute it. The work claims
that disagreement \emph{can be given a measured, auditable form that
sustains itself across domains} --- and offers as evidence the
instrument's coherence, its calibration fixtures, its cross-domain
replication, and its defect record. It does \textbf{not} claim
demonstrated deliberative efficacy, because no controlled study exists.
The study that would settle the question is specifiable and is proposed
here as the work's direct continuation: pairs with a live disagreement
on a mapped domain deliberate twice --- once unaided, once with their
hypermap comparison --- with pre- and post-measures of position location
(can each party state the other's stance to the other's satisfaction?),
perceived understanding, and movement; a third arm substitutes a
conventional argument map (the IBIS form) to separate the effect of
\emph{mapping} from the effect of \emph{measurement}. Until that study
is run, every efficacy claim in the dissertation is scoped to what the
evidence carries, and the candidate asks the committee to note that the
scoping is the work's own doing.

\section{Objection 6 --- ``The founding domain is a sex application. Is
this a serious scholarly
object?''}\label{objection-6-the-founding-domain-is-a-sex-application.-is-this-a-serious-scholarly-object}

\textbf{Rebuttal.} Two answers. The scholarly one: human sexuality is a
legitimate and surveyed field of study --- the project's popularity
model is anchored on the published Joyal (2015) and Lehmiller (2018)
surveys --- and a framework for negotiating desire is a framework for a
real and consequential species of disagreement. The structural one: the
founding domain is the \emph{hardest} test the framework could have been
built under, because it forced the consent architecture --- sovereign
profiles, un-votable limits, gated visibility --- to be rigorous from
the first day rather than retrofitted. That the same engine now runs
politics, faith, chemistry, and consciousness under re-skinned
vocabulary is the generality evidence; that it began where honesty about
desire is hardest is a methodological asset the candidate claims without
embarrassment.

\section{Objection 7 --- ``Your central distinction collapses for
multi-pole topics --- on your own
trace.''}\label{objection-7-your-central-distinction-collapses-for-multi-pole-topics-on-your-own-trace.}

\emph{For topics with three or more poles, the crossed and straight
registers return identical numbers, pole topics appear in no match list,
and group depth scores are zero. The communion/conversation distinction
--- the heart of the thesis --- fails exactly where the framework
generalises.}

\textbf{Open item, with a partial rebuttal.} The facts are as stated;
they are in the dissertation because the candidate's own trace put them
there (2026-11-19). The rebuttal covers the two-flank case, which is
neither small nor degenerate: it is the founding domain entire --- 372
topics --- and there the distinction is a theorem, not an aspiration
(Objection 2). The generalisation to n poles is genuinely open, and the
dissertation names the fork honestly: either aversion lives only on the
intensity coordinate's sign, freeing pole signs to carry a crossed
meaning, or pole signs keep their present semantics and the crossed
register for pole topics must be constructed differently. The
candidate's tentative semantics (crossed = communion; straight = an
interesting conversation) are recorded as tentative. A framework whose
generalisation has a located, stateable failure is in a better epistemic
position than one whose generalisation was never traced; the committee
is asked to weigh the trace itself as part of the contribution.

\section{Objection 8 --- ``Much of this was drafted and computed by AI
collaborators. Whose work is being
examined?''}\label{objection-8-much-of-this-was-drafted-and-computed-by-ai-collaborators.-whose-work-is-being-examined}

\textbf{Rebuttal, from the record.} The authorship model is stated on
the paper and in the artefact's own About text: Jason Mark Wrinkle,
principal investigator and first author --- the questions, the rulings,
the notation, and the physics are his; Yolanda Biche et al., second
author --- drafting and computation; the application credits four AI
collaborators by name. The development record substantiates the
division: every load-bearing decision in this dissertation is a dated
\emph{ruling} by the investigator, including rulings that overruled his
collaborators (the retention of the quality gate against a planner's
extrapolation; the rejection of a derived-intensity reading in favour of
the (n+1) model; the striking of a co-authored sample's middle name on
his partner's correction). The candidate does not claim to have typed
the work. He claims to have \emph{judged} it --- and judgement,
documented at this granularity, is examinable in the viva in the
ordinary way.

\section{Objection 9 --- ``Why rhetoric? This belongs, if anywhere, in
computer science or information
science.''}\label{objection-9-why-rhetoric-this-belongs-if-anywhere-in-computer-science-or-information-science.}

\textbf{Rebuttal.} Because the question is a rhetorical question. The
work asks how disagreement can be \emph{located} --- the founding
question of stasis theory --- and what two people may honestly say to
each other about where they stand. Computer science would examine the
software and find the contribution thin; information science would
examine the classifications and find the comparison mathematics
extraneous. Each neighbouring field amputates half the object. Rhetoric
--- specifically argumentation and deliberation studies --- is the field
whose object is the whole of it: discourse between persons who disagree,
and the instruments by which that discourse can be made honest. The
candidate concedes the standing risk of interdisciplinary work --- that
it is judged by the standards of the field it is not in --- and submits
that Chapter One's lineage is the answer: the work has ancestors in this
field, and it is to this field's questions that its claims are
addressed.

\section{Objection 10 --- ``The stasis parallel is a stretch. Hermagoras
classified disputes; you classify
stances.''}\label{objection-10-the-stasis-parallel-is-a-stretch.-hermagoras-classified-disputes-you-classify-stances.}

\textbf{Concession with boundary.} Correct, and stated in Chapter One as
the delta rather than hidden from it. Lineage is not identity: the claim
is descent of a \emph{founding conviction} --- that the first act in a
dispute is locative, not combative --- together with the taxonomic form
itself (an enumerable scheme of the kinds of a contested thing). The
unit-of-analysis shift from disputes to stances, and from a rhetor's
diagnostic art to a self-administered instrument, is the contribution,
and it is not strengthened by pretending Hermagoras anticipated it.

\section{Objection 11 --- ``The popularity model mixes survey data with
FetLife counts and `literature priors.' Is this measurement or
numerology?''}\label{objection-11-the-popularity-model-mixes-survey-data-with-fetlife-counts-and-literature-priors.-is-this-measurement-or-numerology}

\textbf{Concession with boundary.} The model's sources, weights, and
aggregation rule are stated in full in the companion paper (\S 4) ---
published surveys anchored at factor 1.0; platform counts discounted and
log-normalised; priors weakest and explicitly provisional; missing data
ignored rather than zeroed --- and two of its components (Reddit counts;
the shipped sort key's fidelity to the aggregate) are flagged in the
text as pending and unverified respectively. The committee may
reasonably judge the model's \emph{inputs} thin; it cannot judge them
hidden. That distinction --- between a weak instrument openly specified
and a strong claim opaquely made --- is the dissertation's recurring
standard, applied here to the candidate's own work.

\section{Objection 12 --- ``Chapter Ten of the paper speculates about
estrus cycles and phases of desire. What is physiology doing in a
rhetoric
thesis?''}\label{objection-12-chapter-ten-of-the-paper-speculates-about-estrus-cycles-and-phases-of-desire.-what-is-physiology-doing-in-a-rhetoric-thesis}

\textbf{Concession with boundary.} The penumbra material (Heat as a
phase; the dark sums K and $\epsilon$) is marked provisional in the paper itself,
and the candidate's defense of it is narrow: the \emph{dark sums} are
arithmetic facts about the product that any complete account must either
read or consciously leave unread, and Initiative --- the attribution of
a scene's driving force to the party whose wanting exceeds the appetite
it meets --- is a genuine rhetorical quantity (it distinguishes
authority from initiative, two facts deliberation constantly confuses).
The physiological phase model is an open question the paper poses and
declines to settle. If the committee requires it, the phase material can
be moved to a speculative appendix without touching any defended claim;
the candidate would regard that as a fair price and says so in advance.

\section{What would change the candidate's
mind}\label{what-would-change-the-candidates-mind}

A defense should state its falsifiers. The framework's core claims would
fall, and the candidate would say so, if: (i) the proposed deliberation
study (Objection 5) showed mapped pairs locating each other's positions
\emph{no better} than unaided pairs --- the instrument would then be a
horoscope with receipts, receipts notwithstanding; (ii) a rival skin,
authored in good faith under the same rules, proved inexpressible in the
pole formalism for a domain the framework claims to cover --- the
stance-space would then be a local, not a general, form; or (iii) the
crossed/straight distinction for two-flank topics were shown to make no
difference any user can detect or use --- the two registers would then
be one register with two names, and the thesis's heart would be a
distinction without a difference. None of these has occurred. All three
are runnable. That is the candidate's closing submission: the work is
examinable because it can lose.

\begin{center}\rule{0.5\linewidth}{0.5pt}\end{center}

\emph{End of the three commissioned chapters. The companion tensor paper
(``How Maps Multiply'') is updated in parallel: revision of 2026-11-21.}


\end{document}
